\documentclass[11pt]{amsart}
\usepackage[T1]{fontenc}
\usepackage{lmodern}
\usepackage{amssymb,amsthm}
\usepackage{booktabs}
\usepackage{array}
\usepackage{hyperref}

% ------------------------------------------------------------------
% Theorem environments
% ------------------------------------------------------------------
\theoremstyle{plain}
\newtheorem*{mainthm}{Theorem}
\newtheorem{lemma}{Lemma}
\newtheorem{corollary}[lemma]{Corollary}
\newtheorem*{lemmaG}{Lemma G}
\theoremstyle{definition}
\newtheorem*{conjecture}{Conjecture}
\theoremstyle{remark}
\newtheorem{remark}{Remark}

\numberwithin{equation}{section}

% ------------------------------------------------------------------
% Macros
% ------------------------------------------------------------------
\newcommand{\QG}{\mathcal{Q}}          % the n-Queens graph
\newcommand{\Adj}{\mathbf{A}}          % its adjacency matrix
\newcommand{\R}{\mathbb{R}}
\newcommand{\Z}{\mathbb{Z}}
\newcommand{\iv}[1]{[#1]}              % Iverson bracket
\newcommand{\abs}[1]{\lvert #1\rvert}
\newcommand{\rev}[1]{#1^{\leftarrow}}  % reversal of a row vector
\newcommand{\refl}[1]{#1^{-}}          % reflection of a diagonal vector
\newcommand{\dc}{d^{\circ}}
\newcommand{\Eq}{\mathrm{Eq}}
\newcommand{\VA}{V_{A_1}}
\newcommand{\VBa}{V_{B_1}}
\newcommand{\VBb}{V_{B_2}}
\newcommand{\VE}{V_{E}}
\newcommand{\VEp}{V_{E'}}
\DeclareMathOperator{\rank}{rank}
\newcommand{\leanus}{\textunderscore\allowbreak}
\newcommand{\lean}[1]{{\ttfamily\let\_\leanus #1}}  % a Lean name; may break after "_"
\emergencystretch=2em

\title[The multiplicity of $n-4$ in the $n$-Queens graph]{The exact multiplicity of the eigenvalue $n-4$ of the $n$-Queens graph}
\author{Sukhraj Sachdeva}
\date{Version 3, 2026-11-27. All versions: DOI \href{https://doi.org/10.5281/zenodo.24596852}{10.5281/zenodo.24596852}}

\begin{document}

\begin{abstract}
For every integer \(n\ge 4\) we show that the eigenvalue \(n-4\) of the adjacency
matrix of the \(n\)-Queens graph has multiplicity exactly \((n-2)/2\) when \(n\) is even
and \((n+1)/2\) when \(n\) is odd. Cardoso, Costa and Duarte proved that these numbers
are lower bounds, with explicit eigenvectors, and conjectured equality. We prove the
matching upper bound. The proof is by hand and uses only elementary linear algebra:
three symmetries of the board split the problem into four small linear systems, and
each system is solved explicitly. No step uses a computer. The whole theorem, for every
\(n\ge4\), has also been checked by the Lean~4 proof assistant with the Mathlib library,
as has the explicit description of the eigenspace that the proof yields.
Two other parts of the conjecture of Cardoso, Costa and Duarte remain open.

Human-directed research produced with the help of AI tools (Claude) and checked by AI
referees (Claude Opus, Claude Fable and Claude Sonnet; ChatGPT and Gemini as outside
cross-checks). The main theorem is formalized in Lean~4 with Mathlib. Additional
verification may still be required.
\end{abstract}

\maketitle

% ==================================================================
\section{Introduction}\label{sec:intro}
% ==================================================================

\subsection{The problem}
The \emph{\(n\)-Queens graph} \(\QG(n)\) has the \(n^2\) squares of an \(n\times n\)
chessboard as its vertices. Two distinct squares are adjacent when a chess queen on one
of them attacks the other, that is, when they share a row, a column or a diagonal (in
either direction). Its \emph{adjacency matrix} \(\Adj\) is the \(n^2\times n^2\) matrix
with entry \(1\) for adjacent squares and \(0\) otherwise. It is real and symmetric, so
all its eigenvalues are real. The \emph{multiplicity} of an eigenvalue \(\lambda\) is its
multiplicity as a root of the characteristic polynomial; for a symmetric matrix this
equals \(\dim\ker(\Adj-\lambda I)\).

Cardoso, Costa and Duarte studied which eigenvalues of \(\Adj\) are integers
\cite{CCD20,CCD23,CCD25}. For \(n\ge4\) let \(L(n)\) be the following set of integers:
\(L(n)=\{-4,\,n-4\}\) if \(n\) is even, and
\[
L(n)=\Bigl\{-4,-3,\dots,\tfrac{n-11}{2}\Bigr\}\cup\Bigl\{\tfrac{n-5}{2},\dots,n-5,\,n-4\Bigr\}
\qquad\text{if \(n\) is odd.}
\]
They proved that every element of \(L(n)\) other than \(-4\) is an eigenvalue of \(\Adj\)
\cite[Sections~3--4]{CCD23}, that \(-4\) is the least eigenvalue and has multiplicity
exactly \((n-3)^2\) \cite[Theorems~3.6 and~3.8]{CCD20}, \cite{CCD22}, and that \(n-4\) has
multiplicity at least \((n-2)/2\) (\(n\) even) or \((n+1)/2\) (\(n\) odd), with explicit
eigenvectors. The last result first appeared in \cite[Theorem~3.10]{CCD20} and was
restated in \cite[Theorem~4.1 and Corollary~4.2]{CCD23}. They stated the rest as a
conjecture. Their Conjecture~1 \cite{CCD23} consists of three sentences, which the source
does not number. We restate it in our words and label its three sentences (a), (b) and
(c), in order; these labels are ours. So part~(b) below is the second sentence of
Conjecture~1 of \cite{CCD23}.

\begin{conjecture}[Cardoso, Costa and Duarte {\cite[Conjecture~1]{CCD23}}]
Let \(n\ge4\).
\begin{enumerate}
\item[(a)] The integer eigenvalues of \(\Adj\) are exactly the elements of \(L(n)\).
\item[(b)] The eigenvalue \(n-4\) has multiplicity \((n-2)/2\) if \(n\) is even and
\((n+1)/2\) if \(n\) is odd.
\item[(c)] The elements of \(L(n)\) other than \(-4\) and \(n-4\) are simple eigenvalues.
\end{enumerate}
\end{conjecture}

The conjecture first appeared in \cite[Conjecture~3]{CCD20}. It is restated as open in
\cite[Conjecture~4.1]{CCD25} and in Costa's thesis \cite{Cos24}. The authors report that
their computations confirm it for \(n\le100\) \cite{CCD20,CCD23}; the method is not
stated. In \cite{CCD25} they decompose \(\QG(n)\) into edge-disjoint pieces, among them
\emph{integral} graphs (graphs whose eigenvalues are all integers), and bound
individual eigenvalues with Weyl's inequalities. Our
literature search also covered the spectra of simplicial rook graphs
\cite{MW15,BCHV16}, another family of integral graphs; neither of these papers studies
the \(n\)-Queens graph. Our proof
uses neither them nor the decomposition of \cite{CCD25}.

\subsection{Result}
This paper proves part~(b), the second sentence of Conjecture~1 of \cite{CCD23}.

\begin{mainthm}
Let \(n\ge4\). Then \(n-4\) is an eigenvalue of the adjacency matrix of \(\QG(n)\), and
its multiplicity is exactly \((n-2)/2\) if \(n\) is even and \((n+1)/2\) if \(n\) is odd.
\end{mainthm}

Since the lower bounds are known \cite[Theorem~3.10]{CCD20}, \cite[Corollary~4.2]{CCD23},
the content of the Theorem is the upper bound. Our proof does not use the lower bound: it computes the dimension of
the eigenspace exactly. As a by-product it describes the eigenspace
(Corollary~\ref{cor:eigenspace}); that description coincides with the span of the
eigenvectors of Cardoso, Costa and Duarte (Remark~\ref{rem:credit}).

Parts (a) and (c) of the conjecture remain open. This paper does not settle the
conjecture.

\subsection{Credit}
The conjecture, the lower bounds for the multiplicity of \(n-4\), and the eigenvectors
are due to Cardoso, Costa and Duarte \cite{CCD20,CCD23}. The description of the
eigenspace in Corollary~\ref{cor:eigenspace} is theirs in full: their eigenvectors span
it. The exact value, equivalently the upper bound (the Theorem), is conjectured but not proved in their papers (\cite[Conjecture~3]{CCD20}; \cite[Conjecture~1]{CCD23}; \cite[Conjecture~4.1]{CCD25}); whether a proof of it,
or the method used here, appears elsewhere has not been checked beyond our literature search of 2026-11-27.
That search found no earlier proof or
claim of this upper bound in the sources we checked. Some sources could not be checked that day; they
are listed in the file \nolinkurl{LITERATURE_CHECK.md} of the accompanying repository:
\begin{center}
\footnotesize\url{https://github.com/nebulousice/cardoso-costa-duarte-queens-conjecture}
\end{center}
This is a dated search result, not a priority claim.

\subsection{Method}
The proof has four steps.
\begin{enumerate}
\item Every square lies on exactly four \emph{lines} (its row, column, diagonal and
antidiagonal), and two distinct squares are adjacent exactly when they share a line.
So \(\Adj=NN^{T}-4I\), where \(N\) is the square--line incidence matrix. It follows that
the multiplicity of \(n-4\) in \(\Adj\) equals \(\dim\ker(M-nI)\), where \(M=N^{T}N\) is
the \((6n-2)\times(6n-2)\) matrix of intersection sizes of lines
(Section~\ref{sec:reduction}).
\item Three explicit symmetries of the board (transpose, left--right flip, half-turn)
split the space of functions on lines into five pieces that \(M\) preserves. Two of the
pieces behave identically (Section~\ref{sec:sectors}).
\item On each piece, \(M\) has an explicit formula (Section~\ref{sec:action}), and the
kernel of \(M-nI\) is found by hand: dimension \(\lceil n/2\rceil-1\) on one piece,
\(1\) or \(0\) on another depending on the parity of \(n\), and \(0\) on the other three
(Sections~\ref{sec:B1E}--\ref{sec:B2}). The hardest piece is handled by a discrete
maximum principle (Lemma~\ref{lem:means}).
\item Adding up gives the Theorem (Section~\ref{sec:main}).
\end{enumerate}
No representation theory is needed: the decomposition comes from three explicit
involutions.

\subsection{What is proved by hand, what is checked by computer, and what is
formalized}
\begin{itemize}
\item \emph{By hand:} everything. The Theorem, every lemma and Corollary~\ref{cor:eigenspace}
have complete written proofs below. No step uses a computation.
\item \emph{By computer, as corroboration only:} the Theorem and its individual steps,
for finitely many \(n\), in exact integer and rational arithmetic
(Section~\ref{sec:computer}). These checks are not certificates in the sense of
Section~\ref{sec:computer}, and the proof does not use them.
\item \emph{Formalized in a proof assistant:} the whole Theorem, for every \(n\ge4\),
in Lean~4 with Mathlib (Section~\ref{sec:formal}). The Lean proof follows the
same plan as the written proof, with a few local differences listed there.
Corollary~\ref{cor:eigenspace} is also formalized, as \lean{queens\_eigenspace} (over
any field with a compatible total order, such as \(\mathbb{Q}\) or \(\R\)) and \lean{queens\_eigenspace\_real}, in the file
\nolinkurl{Eigenspace.lean}. The appendix is not formalized.
\end{itemize}
The written proof does not depend on the computer checks or on the Lean files. The Lean
files give a separate, machine-checked proof of the same Theorem, starting from the
definition of the graph.

\subsection{Status of this work}
The proof was written on 2026-11-27. It has been through three review rounds by AI
referees (Claude models), each working independently and checking the proof line by line
and with their own exact computations. No referee has found a mathematical error. In the
first round, one referee accepted with no issues; the other raised three minor issues
(an unsupported status label for the finite-range computations, incomplete credit for
the eigenspace, and a mislabelled comment and a stale sentence), all fixed without
changing any mathematical step. In the second round, a third referee accepted with no
issues; the other two raised minor issues about wording and citation (a stale status
sentence, and that the lower bound should be credited first to \cite{CCD20}). In the third round, on the
fixed text, two referees accepted with no issues and the third raised two minor wording
points in the credit remarks (what the papers of Cardoso, Costa and Duarte state about the
upper bound, and a linear relation among the vectors of the printed set \(Y^n\) of
\cite{CCD20}); both are fixed here, with no mathematical change. The proof therefore passed
the review rule we applied: at least two of three referees accept with no issues, and no
referee has an unresolved major issue. The referees read an earlier presentation of the
same proof, with the same lemmas and the same mathematical steps; the present text
reorders the parts of one lemma and rewrites the introductory and closing sections.

After the review, the Theorem was formalized in Lean~4 with Mathlib, for every
\(n\ge4\) (Section~\ref{sec:formal}). The Lean check passed in two separate runs. A
further AI referee (Claude Opus) rebuilt the Lean files from a fresh copy and checked
that the Lean statement says what the source conjecture says. Corollary~\ref{cor:eigenspace}
was then formalized as well, and the same model checked its Lean statement. Finally the
compiled Lean proofs were checked again by several proof-checking programs (kernels),
two of them written independently of Lean (Section~\ref{sec:formal}).

Two other AI systems, ChatGPT and Gemini, were each given the first version of this
paper in a fresh session, with a referee prompt asking for a line-by-line check. They
found no mathematical error and reported minor formatting issues only.

\textbf{Expert verification required.} Verified so far with: Lean~4 and Mathlib (the
whole Theorem and Corollary~\ref{cor:eigenspace}, for every \(n\ge4\); no unproved step;
only the three standard axioms); replays of the compiled Lean proofs by Lean's own
checker and by the independent kernels \texttt{nanoda\_bin} and \texttt{lean4lean}, all
passing; three rounds of AI review of the written proof by Claude Opus, Claude Fable and
Claude Sonnet, as described above (passed in round~3: Claude Opus and Claude Sonnet
accepted with no issues, Claude Fable raised two minor wording points, since fixed; no
referee found a mathematical error or raised a major issue); AI checks of the Lean
statements, of the Theorem against the source conjecture and of
Corollary~\ref{cor:eigenspace} against its written statement (Claude Opus; both
confirmed); exact computation for
\(4\le n\le40\) (Section~\ref{sec:computer}), which is corroboration and not a
certificate; and the outside cross-checks by ChatGPT and Gemini described above.

% ==================================================================
\section{Notation and conventions}\label{sec:notation}
% ==================================================================

Throughout, \(n\ge4\) is an integer, and for a statement \(P\) we write \(\iv{P}=1\) if
\(P\) holds and \(\iv{P}=0\) otherwise.

\emph{Board.} The squares are the pairs \((i,j)\) with \(0\le i,j\le n-1\). Two distinct
squares \((i,j)\) and \((i',j')\) are adjacent in \(\QG(n)\) when \(i=i'\), or \(j=j'\),
or \(i-j=i'-j'\), or \(i+j=i'+j'\). (The sources index squares from \(1\); shifting the
indices is a graph isomorphism and changes nothing.)

\emph{Lines.} There are four families of lines: the rows \(R_i\) and columns \(C_j\)
(\(0\le i,j\le n-1\)), the diagonals \(D_k\) and the antidiagonals \(A_s\)
(\(\abs{k},\abs{s}\le n-1\)):
\begin{align*}
R_i&=\{(i,j)\mid 0\le j\le n-1\}, & C_j&=\{(i,j)\mid 0\le i\le n-1\},\\
D_k&=\{(i,j)\mid i-j=k\}, & A_s&=\{(i,j)\mid i+j-(n-1)=s\}.
\end{align*}
There are \(n+n+(2n-1)+(2n-1)=6n-2\) lines, all nonempty. The square \((i,j)\) lies on
exactly four lines, one of each family: \(R_i\), \(C_j\), \(D_{i-j}\) and
\(A_{i+j-(n-1)}\). (The letter \(A\) with an integer subscript always denotes an
antidiagonal; the adjacency matrix is the bold \(\Adj\).)

\emph{Line space.} Let \(V=\R^{6n-2}\) be the space of real functions on the set of
lines. We write a vector as \(v=(r,c,d,a)\), where \(r=(r_i)_{0\le i\le n-1}\) lists its
values on the rows, \(c=(c_j)_{0\le j\le n-1}\) on the columns,
\(d=(d_k)_{\abs{k}\le n-1}\) on the diagonals and \(a=(a_s)_{\abs{s}\le n-1}\) on the
antidiagonals.

\emph{Incidence and Gram matrices.} Let \(N\) be the \(n^2\times(6n-2)\) matrix with
entries \(N_{x,L}=\iv{x\in L}\) for a square \(x\) and a line \(L\), and let
\(M=N^{T}N\). Then \(M_{L,L'}=\abs{L\cap L'}\). For \(v=(r,c,d,a)\),
\begin{equation}\label{eq:Nv}
(Nv)(i,j)=r_i+c_j+d_{i-j}+a_{i+j-(n-1)},\qquad (Mv)_L=\sum_{x\in L}(Nv)(x).
\end{equation}

\emph{Reversals.} For \(r\in\R^n\) put \(\rev{r}_i=r_{n-1-i}\). For
\(d=(d_k)_{\abs{k}\le n-1}\) put \(\refl{d}_k=d_{-k}\). We call \(r\)
\emph{palindromic} if \(\rev{r}=r\) and \emph{antipalindromic} if \(\rev{r}=-r\). We call
\(d\) \emph{even} if \(\refl{d}=d\) and \emph{odd} if \(\refl{d}=-d\); an odd \(d\) has
\(d_0=0\).

\emph{Sums.} For \(r\in\R^n\) put \(\rho_k=\sum_{i=k}^{n-1}r_i\) for \(0\le k\le n-1\),
\(\rho_n=0\) and \(R=\rho_0=\sum_i r_i\). Then \(r_i=\rho_i-\rho_{i+1}\). For \(m\ge0\) let
\(S(m)=\{-m,-m+2,\dots,m-2,m\}\); it has \(m+1\) elements and is symmetric under
\(s\mapsto -s\). For real numbers \(d_0,\dots,d_{n-1}\) put
\begin{equation}\label{eq:T}
T_m=\sum_{s\in S(m)}d_{\abs{s}}\quad(0\le m\le n-1),\qquad T_{-1}=0,\qquad
D=\sum_{k=0}^{n-1}d_k.
\end{equation}
Thus \(T_m=d_0+2(d_2+d_4+\dots+d_m)\) for even \(m\), and \(T_m=2(d_1+d_3+\dots+d_m)\) for
odd \(m\). (Without a subscript, \(R\) and \(D\) denote the sums just defined; with a
subscript, \(R_i\) and \(D_k\) denote lines.)

% ==================================================================
\section{Reduction to the Gram matrix of the lines}\label{sec:reduction}
% ==================================================================

This section shows that the multiplicity of \(n-4\) in \(\Adj\) equals
\(\dim\ker(M-nI)\). Both lemmas are standard facts about edge clique partitions and Gram
matrices (compare \cite{CCD22}); we prove them to keep the paper self-contained.

\begin{lemma}\label{lem:clique}
Two distinct squares lie on at most one common line, and they are adjacent in \(\QG(n)\)
if and only if they lie on a common line. Hence \(NN^{T}=\Adj+4I\).
\end{lemma}

\begin{proof}
The entry \((NN^{T})_{x,x'}\) is the number of lines that contain both \(x\) and \(x'\).
For \(x=x'\) it is \(4\), since every square lies on exactly four lines. Let
\(x=(i,j)\ne x'=(i',j')\). Each square lies on exactly one line of each family, so the
number of common lines is the number of true statements among
\[
\text{(a) } i=i',\qquad \text{(b) } j=j',\qquad \text{(c) } i-j=i'-j',\qquad
\text{(d) } i+j=i'+j'.
\]
By definition, \(x\) and \(x'\) are adjacent exactly when at least one of them holds. Any
two of them force \(x=x'\): (a) and (b) directly; (a) and (c) give \(j=j'\); (a) and (d)
give \(j=j'\); (b) and (c) give \(i=i'\); (b) and (d) give \(i=i'\); (c) and (d) give
\(2i=2i'\) by addition, hence \(i=i'\) and then \(j=j'\). So for \(x\ne x'\) at most one
statement holds, and \((NN^{T})_{x,x'}=\iv{x\text{ adjacent to }x'}=\Adj_{x,x'}\).
\end{proof}

\begin{lemma}\label{lem:gram}
For every real \(\mu\ne0\), the map \(v\mapsto Nv\) is a linear bijection from
\(\ker(M-\mu I)\) onto \(\ker(NN^{T}-\mu I)\). Consequently, the multiplicity of \(n-4\)
as an eigenvalue of \(\Adj\) equals \(\dim\ker(M-nI)\), and \(N\) maps \(\ker(M-nI)\)
bijectively onto \(\ker(\Adj-(n-4)I)\).
\end{lemma}

\begin{proof}
If \(Mv=\mu v\), then \(NN^{T}(Nv)=N(Mv)=\mu Nv\). So \(N\) maps \(\ker(M-\mu I)\) into
\(\ker(NN^{T}-\mu I)\). It is injective there: if \(Mv=\mu v\) and \(Nv=0\), then
\(\mu v=N^{T}Nv=0\), so \(v=0\) because \(\mu\ne0\). Hence
\(\dim\ker(M-\mu I)\le\dim\ker(NN^{T}-\mu I)\). The same argument with \(N\) and \(N^{T}\)
exchanged gives the reverse inequality. So the dimensions are equal, and the injective
map \(N\) is a bijection. By Lemma~\ref{lem:clique},
\(\ker(\Adj-(n-4)I)=\ker(NN^{T}-nI)\), and \(\mu=n\ne0\).
\end{proof}

% ==================================================================
\section{Three symmetries and the five sectors}\label{sec:sectors}
% ==================================================================

This section splits \(V\) into five subspaces that \(M\) preserves. We use three
explicit \emph{involutions} (linear maps whose square is the identity). The names
\(A_1\), \(B_1\), \(B_2\), \(E\), \(E'\) of the subspaces are only labels.

Define linear maps \(t,h,z\colon V\to V\) by
\begin{equation}\label{eq:ops}
\begin{gathered}
t(r,c,d,a)=(c,r,\refl{d},a),\qquad h(r,c,d,a)=(r,\rev{c},a,d),\\
z(r,c,d,a)=(\rev{r},\rev{c},\refl{d},\refl{a}).
\end{gathered}
\end{equation}

\begin{lemma}\label{lem:sym}
\begin{enumerate}
\item[(a)] \(t^2=h^2=z^2=I\), \(tz=zt\), \(hz=zh\), \(th=zht\), and \((th)^2=z\).
\item[(b)] Each of \(t\), \(h\), \(z\) commutes with \(M\).
\end{enumerate}
\end{lemma}

\begin{proof}
(a) The reversals \(r\mapsto\rev{r}\) and \(d\mapsto\refl{d}\) are involutions, and so is
a swap of two entries of \((r,c,d,a)\); hence \(t^2=h^2=z^2=I\). For \(v=(r,c,d,a)\),
substituting into \eqref{eq:ops} gives
\begin{align*}
tzv&=t(\rev{r},\rev{c},\refl{d},\refl{a})=(\rev{c},\rev{r},d,\refl{a}),\\
ztv&=z(c,r,\refl{d},a)=(\rev{c},\rev{r},d,\refl{a}),\\
hzv&=h(\rev{r},\rev{c},\refl{d},\refl{a})=(\rev{r},c,\refl{a},\refl{d}),\\
zhv&=z(r,\rev{c},a,d)=(\rev{r},c,\refl{a},\refl{d}),\\
thv&=t(r,\rev{c},a,d)=(\rev{c},r,\refl{a},d),\\
zhtv&=z(c,\rev{r},a,\refl{d})=(\rev{c},r,\refl{a},d).
\end{align*}
So \(tz=zt\), \(hz=zh\) and \(th=zht\). Finally \((th)^2=(zht)(th)=zh(tt)h=zhh=z\).

(b) Consider the three bijections of the board \(\tau(i,j)=(j,i)\) (transpose),
\(\eta(i,j)=(i,n-1-j)\) (left--right flip) and \(\zeta(i,j)=(n-1-i,n-1-j)\) (half-turn).
Each is an involution, and each maps lines onto lines as in the following table.
\begin{center}
\begin{tabular}{@{}lcccc@{}}
\toprule
line \(L\) & \(R_i\) & \(C_j\) & \(D_k\) & \(A_s\)\\
\midrule
\(\tau(L)\) & \(C_i\) & \(R_j\) & \(D_{-k}\) & \(A_s\)\\
\(\eta(L)\) & \(R_i\) & \(C_{n-1-j}\) & \(A_k\) & \(D_s\)\\
\(\zeta(L)\) & \(R_{n-1-i}\) & \(C_{n-1-j}\) & \(D_{-k}\) & \(A_{-s}\)\\
\bottomrule
\end{tabular}
\end{center}
Each entry \(\sigma(L)\subseteq L'\) is checked by substituting into the defining
equation of the line. For instance, \((i,j)\in D_k\) means \(i-j=k\), which is
\(i+(n-1-j)-(n-1)=k\); so \(\eta(i,j)\in A_k\). The inclusions are equalities: the table
also gives \(\sigma(L')\subseteq L\) (for example \(\eta(A_k)\subseteq D_k\)), hence
\(L'=\sigma(\sigma(L'))\subseteq\sigma(L)\).

Let \(\sigma\) be one of \(\tau,\eta,\zeta\). It permutes the lines, and
\(\sigma^{-1}=\sigma\). Define \(\hat\sigma\colon V\to V\) by
\((\hat\sigma v)_L=v_{\sigma(L)}\). Since \(\sigma\) is a bijection of the board,
\(\abs{\sigma(L)\cap\sigma(L')}=\abs{L\cap L'}\) for all lines \(L,L'\). Substituting
\(L'=\sigma(L'')\) and then applying \(\sigma\) to both sets,
\begin{align*}
(M\hat\sigma v)_L&=\sum_{L'}\abs{L\cap L'}\,v_{\sigma(L')}
=\sum_{L''}\abs{L\cap\sigma(L'')}\,v_{L''}\\
&=\sum_{L''}\abs{\sigma(L)\cap L''}\,v_{L''}=(Mv)_{\sigma(L)}=(\hat\sigma Mv)_L.
\end{align*}
So \(\hat\sigma\) commutes with \(M\). Reading the table,
\[
\hat\tau v=(c,r,\refl{d},a),\qquad
\hat\eta v=(r,\rev{c},a,d),\qquad
\hat\zeta v=(\rev{r},\rev{c},\refl{d},\refl{a});
\]
that is, \(\hat\tau=t\), \(\hat\eta=h\) and \(\hat\zeta=z\).
\end{proof}

Define six subspaces of \(V\):
\begin{align*}
\VA&=\{v\in V\mid tv=v,\ hv=v\}, & \VBa&=\{v\in V\mid tv=-v,\ hv=v\},\\
\VBb&=\{v\in V\mid tv=v,\ hv=-v\}, & W&=\{v\in V\mid tv=-v,\ hv=-v\},\\
\VE&=\{v\in V\mid zv=-v,\ hv=v\}, & \VEp&=\{v\in V\mid zv=-v,\ hv=-v\}.
\end{align*}
For a subspace \(U\) with \(MU\subseteq U\) and a real \(\mu\), write
\(\kappa_U(\mu)=\dim\ker\bigl((M-\mu I)|_U\bigr)\).

\begin{lemma}\label{lem:sectors}
\begin{enumerate}
\item[(a)] \(W=0\), and
\begin{align}
\VA&=\{(r,r,d,d)\mid r\text{ palindromic},\ d\text{ even}\},\label{eq:VA}\\
\VBa&=\{(r,-r,0,0)\mid r\text{ palindromic}\},\label{eq:VB1}\\
\VBb&=\{(0,0,d,-d)\mid d\text{ even}\},\label{eq:VB2}\\
\VE&=\{(r,0,d,d)\mid r\text{ antipalindromic},\ d\text{ odd}\}.\label{eq:VE}
\end{align}
Their dimensions are \(\lceil n/2\rceil+n\), \(\lceil n/2\rceil\), \(n\) and
\(\lfloor n/2\rfloor+n-1\).
\item[(b)] \(V=\VA\oplus\VBa\oplus\VBb\oplus\VE\oplus\VEp\).
\item[(c)] Each of the five subspaces in (b) is \(M\)-invariant, and \(t\) restricts to
a linear bijection \(\VE\to\VEp\) that commutes with \(M\).
\item[(d)] For every real \(\mu\),
\[
\dim\ker(M-\mu I)=\kappa_{\VA}(\mu)+\kappa_{\VBa}(\mu)+\kappa_{\VBb}(\mu)
+2\kappa_{\VE}(\mu).
\]
\end{enumerate}
\end{lemma}

\begin{proof}
Part (a) uses only \eqref{eq:ops}. Part (b) uses Lemma~\ref{lem:sym}(a) and part (a).
Part (c) uses Lemma~\ref{lem:sym}. Part (d) uses parts (b) and (c).

(a) Write \(v=(r,c,d,a)\). By \eqref{eq:ops}:
\(tv=v\) means \(c=r\) and \(d\) even;
\(tv=-v\) means \(c=-r\), \(d\) odd and \(a=0\);
\(hv=v\) means \(c\) palindromic and \(a=d\);
\(hv=-v\) means \(r=0\), \(c\) antipalindromic and \(a=-d\);
\(zv=-v\) means \(r\) and \(c\) antipalindromic and \(d\) and \(a\) odd.
Combining these: for \(\VA\), \(c=r\) is palindromic and \(a=d\) is even, which is
\eqref{eq:VA}. For \(\VBa\), \(c=-r\) is palindromic and \(d=a=0\), which is
\eqref{eq:VB1}. For \(\VBb\), \(r=0\), \(c=r=0\), \(d\) is even and \(a=-d\), which is
\eqref{eq:VB2}. For \(W\), \(r=0\), \(c=-r=0\), \(a=0\) and \(d=-a=0\); so \(W=0\). For
\(\VE\), \(c\) is both palindromic and antipalindromic, so \(c=0\); \(r\) is
antipalindromic, and \(a=d\) is odd; this is \eqref{eq:VE}. The dimensions follow by
counting free coordinates: a palindromic \(r\) has \(\lceil n/2\rceil\) free entries, an
antipalindromic one \(\lfloor n/2\rfloor\), an even \(d\) has \(n\) (namely
\(d_0,\dots,d_{n-1}\)), and an odd \(d\) has \(n-1\) (namely \(d_1,\dots,d_{n-1}\)).

(b) Put \(V_{\pm}=\ker(z\mp I)\). Since \(z^2=I\), every \(v\) equals
\(\frac12(v+zv)+\frac12(v-zv)\), with the first term in \(V_+\) and the second in
\(V_-\), and \(V_+\cap V_-=0\). So \(V=V_+\oplus V_-\). Since \(t\) and \(h\) commute with
\(z\) (Lemma~\ref{lem:sym}(a)), they map \(V_+\) into \(V_+\) and \(V_-\) into \(V_-\).

\emph{The part \(V_+\).} For \(v\in V_+\) we have \(htv\in V_+\), hence
\(thv=z(htv)=htv\) by Lemma~\ref{lem:sym}(a). So \(t\) and \(h\) restrict to commuting
involutions of \(V_+\). For signs \(\varepsilon,\delta\in\{1,-1\}\) put
\(\pi_{\varepsilon\delta}=\frac14(I+\varepsilon t)(I+\delta h)\) on \(V_+\). Then
\(\sum_{\varepsilon,\delta}\pi_{\varepsilon\delta}=\frac14(2I)(2I)=I\). Since \(t\) and
\(h\) commute on \(V_+\), and \(t(I+\varepsilon t)=\varepsilon(I+\varepsilon t)\),
\(h(I+\delta h)=\delta(I+\delta h)\), every vector \(w=\pi_{\varepsilon\delta}v\) with
\(v\in V_+\) satisfies \(tw=\varepsilon w\) and \(hw=\delta w\). Conversely, if
\(w\in V_+\), \(tw=\varepsilon'w\) and \(hw=\delta'w\), then
\(\pi_{\varepsilon\delta}w=\frac14(1+\varepsilon\varepsilon')(1+\delta\delta')w\), which
is \(w\) if \((\varepsilon,\delta)=(\varepsilon',\delta')\) and \(0\) otherwise. Hence
\(V_+\) is the direct sum of the four spaces
\(\{v\in V_+\mid tv=\varepsilon v,\ hv=\delta v\}\): they span \(V_+\), and applying
\(\pi_{\varepsilon\delta}\) to a vanishing sum of vectors, one from each space, shows
that each summand is zero. Moreover, any \(v\in V\) with \(tv=\varepsilon v\) and
\(hv=\delta v\) lies in \(V_+\), because \(zv=(th)^2v=(\varepsilon\delta)^2v=v\). So the
four spaces are \(\VA\), \(\VBa\), \(\VBb\) and \(W\). By part (a), \(W=0\). Hence
\(V_+=\VA\oplus\VBa\oplus\VBb\).

\emph{The part \(V_-\).} Since \(h\) is an involution of \(V_-\), the same argument as
for \(z\) gives \(V_-=\{v\in V_-\mid hv=v\}\oplus\{v\in V_-\mid hv=-v\}=\VE\oplus\VEp\).
This proves (b).

(c) If \(g\in\{t,h,z\}\) and \(gv=\varepsilon v\), then \(g(Mv)=M(gv)=\varepsilon Mv\) by
Lemma~\ref{lem:sym}(b). Each of the five subspaces is defined by conditions of this
form, so it is \(M\)-invariant. Let \(v\in\VE\). Then \(z(tv)=t(zv)=-tv\). Multiplying
\(th=zht\) on the left by \(z\) gives \(ht=zth\), so \(h(tv)=zt(hv)=z(tv)=-tv\). Hence
\(tv\in\VEp\). In the same way, for \(v\in\VEp\) we get \(z(tv)=-tv\) and
\(h(tv)=zt(hv)=-z(tv)=tv\), so \(tv\in\VE\). Since \(t^2=I\), the map
\(t\colon\VE\to\VEp\) is a bijection with inverse \(t\), and it commutes with \(M\).

(d) Let \(V=\bigoplus_i U_i\) with \(MU_i\subseteq U_i\), and write \(v=\sum_i u_i\) with
\(u_i\in U_i\). Then \((M-\mu I)v=\sum_i(M-\mu I)u_i\) with \((M-\mu I)u_i\in U_i\), so
\((M-\mu I)v=0\) if and only if \((M-\mu I)u_i=0\) for all \(i\). Hence
\(\ker(M-\mu I)=\bigoplus_i\ker\bigl((M-\mu I)|_{U_i}\bigr)\). Apply this to the
decomposition (b), whose summands are \(M\)-invariant by (c). By (c), \(t\) maps
\(\ker((M-\mu I)|_{\VE})\) bijectively onto \(\ker((M-\mu I)|_{\VEp})\), so
\(\kappa_{\VEp}(\mu)=\kappa_{\VE}(\mu)\).
\end{proof}

As a check (not used later), the dimensions in (a), with \(\VE\) counted twice, add up to
\(2\lceil n/2\rceil+2\lfloor n/2\rfloor+4n-2=6n-2=\dim V\).

% ==================================================================
\section{The action of \texorpdfstring{\(M\)}{M} on the sectors}\label{sec:action}
% ==================================================================

\begin{lemma}\label{lem:action}
For \(v=(r,c,d,a)\in V\), \(0\le i\le n-1\) and \(0\le k\le n-1\),
\begin{align}
(Mv)_{R_i}&=nr_i+\sum_{j=0}^{n-1}c_j+\sum_{k'=i-n+1}^{i}d_{k'}
+\sum_{s=i-n+1}^{i}a_s,\label{eq:Mrow}\\
(Mv)_{D_k}&=(n-k)d_k+\sum_{i'=k}^{n-1}r_{i'}+\sum_{j=0}^{n-1-k}c_j
+\sum_{s\in S(n-1-k)}a_s.\label{eq:Mdiag}
\end{align}
Consequently, with the notation \eqref{eq:T}:
\begin{enumerate}
\item[(i)] for \(v=(r,r,d,d)\in\VA\):
\begin{equation}\label{eq:actA}
\begin{gathered}
(Mv)_{R_i}=nr_i+R+2\sum_{k'=i-n+1}^{i}d_{\abs{k'}},\\
(Mv)_{D_k}=(n-k)d_k+2\rho_k+T_{n-1-k};
\end{gathered}
\end{equation}
\item[(ii)] for \(v=(r,-r,0,0)\in\VBa\):
\begin{equation}\label{eq:actB1}
(Mv)_{R_i}=nr_i-R;
\end{equation}
\item[(iii)] for \(v=(0,0,d,-d)\in\VBb\):
\begin{equation}\label{eq:actB2}
(Mv)_{D_k}=(n-k)d_k-T_{n-1-k};
\end{equation}
\item[(iv)] for \(v=(r,0,d,d)\in\VE\) and \(1\le k\le n-1\):
\begin{equation}\label{eq:actE}
(Mv)_{R_i}=nr_i+2\sum_{k'=i-n+1}^{i}d_{k'},\qquad (Mv)_{D_k}=(n-k)d_k+\rho_k.
\end{equation}
\end{enumerate}
\end{lemma}

\begin{proof}
By \eqref{eq:Nv}, \((Mv)_L=\sum_{(i,j)\in L}\bigl(r_i+c_j+d_{i-j}+a_{i+j-(n-1)}\bigr)\).

For \(L=R_i\) the squares are \((i,j)\), \(0\le j\le n-1\). As \(j\) runs through
\(0,\dots,n-1\), the index \(i-j\) runs once through \(i-n+1,\dots,i\), and so does the
index \(i+j-(n-1)\). This is \eqref{eq:Mrow}.

For \(L=D_k\) with \(0\le k\le n-1\) the squares are \((i',i'-k)\) with
\(k\le i'\le n-1\); there are \(n-k\) of them. The row index \(i'\) runs through
\(k,\dots,n-1\). The column index \(i'-k\) runs through \(0,\dots,n-1-k\). The
antidiagonal index is \(i'+(i'-k)-(n-1)=-(n-1-k)+2(i'-k)\), which runs once through
\(S(n-1-k)\). This is \eqref{eq:Mdiag}.

(i) Here \(c=r\), so \(\sum_jc_j=R\). Since \(a=d\) is even, both sums over \(d\) and
\(a\) in \eqref{eq:Mrow} equal \(\sum_{k'=i-n+1}^{i}d_{\abs{k'}}\). In \eqref{eq:Mdiag},
\[
\sum_{j=0}^{n-1-k}c_j=\sum_{j=0}^{n-1-k}r_{n-1-j}=\sum_{i'=k}^{n-1}r_{i'}=\rho_k
\]
because \(r\) is palindromic, and
\(\sum_{s\in S(n-1-k)}a_s=\sum_{s\in S(n-1-k)}d_{\abs{s}}=T_{n-1-k}\).
(ii) Put \(c=-r\) and \(d=a=0\) in \eqref{eq:Mrow}.
(iii) Put \(r=c=0\) and \(a_s=-d_{\abs{s}}\) in \eqref{eq:Mdiag}.
(iv) Put \(c=0\) and \(a=d\) in \eqref{eq:Mrow}. In \eqref{eq:Mdiag},
\(\sum_{s\in S(n-1-k)}a_s=\sum_{s\in S(n-1-k)}d_s=0\), because \(S(n-1-k)\) is symmetric
under \(s\mapsto-s\) and \(d\) is odd.
\end{proof}

\begin{remark}[which components decide]\label{rem:components}
Let \(U\) be one of \(\VA,\VBa,\VBb,\VE\) and let \(v\in U\). By
Lemma~\ref{lem:sectors}(c), \((M-nI)v\in U\). By \eqref{eq:VA}--\eqref{eq:VE}, a vector
of \(\VA\) is determined by its components on the \(R_i\) (\(0\le i\le n-1\)) and the
\(D_k\) (\(0\le k\le n-1\)); a vector of \(\VBa\) by its components on the \(R_i\); a
vector of \(\VBb\) by its components on the \(D_k\) (\(0\le k\le n-1\)); and a vector of
\(\VE\) by its components on the \(R_i\) (\(0\le i\le n-1\)) and the \(D_k\)
(\(1\le k\le n-1\)). So \((M-nI)v=0\) if and only if those components of \((M-nI)v\)
vanish. Lemma~\ref{lem:action} gives exactly those components.
\end{remark}

% ==================================================================
\section{The sectors \texorpdfstring{\(B_1\)}{B1} and \texorpdfstring{\(E\)}{E} at
\texorpdfstring{\(\mu=n\)}{mu=n}}\label{sec:B1E}
% ==================================================================

\begin{lemma}[sector \(B_1\)]\label{lem:B1}
We have
\[
\ker\bigl((M-nI)|_{\VBa}\bigr)=\Bigl\{(r,-r,0,0)\Bigm| r\text{ palindromic},\ \sum_ir_i=0\Bigr\},
\]
and \(\kappa_{\VBa}(n)=\lceil n/2\rceil-1\).
\end{lemma}

\begin{proof}
Let \(v=(r,-r,0,0)\in\VBa\). By \eqref{eq:actB1}, the \(R_i\)-component of \((M-nI)v\) is
\(-R\) for every \(i\). By Remark~\ref{rem:components}, \((M-nI)v=0\) if and only if
\(R=0\). The linear form \(r\mapsto R\) on the \(\lceil n/2\rceil\)-dimensional space of
palindromic vectors is not zero (it takes the value \(n\) at \(r=(1,\dots,1)\)), so its
kernel has dimension \(\lceil n/2\rceil-1\).
\end{proof}

\begin{lemma}[sector \(E\)]\label{lem:E}
\(\kappa_{\VE}(n)=0\).
\end{lemma}

\begin{proof}
We write out the equations, difference consecutive row equations, and use a symmetry of
the partial sums \(\rho_k\). Let \(v=(r,0,d,d)\in\VE\), with \(r\) antipalindromic and
\(d\) odd, and suppose \((M-nI)v=0\). By \eqref{eq:actE} and
Remark~\ref{rem:components} this means
\begin{align*}
(\mathrm{row}_i)&\quad 2\sum_{k'=i-n+1}^{i}d_{k'}=0\quad(0\le i\le n-1),\\
(\mathrm{diag}_k)&\quad -kd_k+\rho_k=0\quad(1\le k\le n-1).
\end{align*}

\emph{Step 1.} For \(1\le i\le n-1\), subtract \((\mathrm{row}_{i-1})\) from
\((\mathrm{row}_i)\). The window of summation moves from \([i-n,i-1]\) to
\([i-n+1,i]\), so the difference is \(2(d_i-d_{i-n})=2(d_i+d_{n-i})\), because \(d\) is
odd and \(1\le n-i\le n-1\). Hence \(d_i+d_{n-i}=0\) for \(1\le i\le n-1\).

\emph{Step 2.} Since \(r\) is antipalindromic, \(\rho_0=\sum_ir_i=0\), and for
\(0\le k\le n\),
\[
\rho_{n-k}=\sum_{i\ge n-k}r_i=-\sum_{i\ge n-k}r_{n-1-i}=-\sum_{j=0}^{k-1}r_j
=-(\rho_0-\rho_k)=\rho_k.
\]

\emph{Step 3.} By \((\mathrm{diag}_k)\), \(d_k=\rho_k/k\) for \(1\le k\le n-1\).
Inserting this and Step~2 into Step~1 gives, for \(1\le k\le n-1\),
\[
0=d_k+d_{n-k}=\frac{\rho_k}{k}+\frac{\rho_{n-k}}{n-k}
=\rho_k\Bigl(\frac1k+\frac1{n-k}\Bigr).
\]
The bracket is positive, so \(\rho_k=0\) for \(1\le k\le n-1\). Together with
\(\rho_0=\rho_n=0\) this gives \(r_i=\rho_i-\rho_{i+1}=0\) for all \(i\), and
\(d_k=\rho_k/k=0\) for \(1\le k\le n-1\). So \(v=0\).
\end{proof}

% ==================================================================
\section{The sector \texorpdfstring{\(A_1\)}{A1} at
\texorpdfstring{\(\mu=n\)}{mu=n}}\label{sec:A1}
% ==================================================================

The key to this sector is an identity for the sums \(T_m\) of \eqref{eq:T}, valid when
\(d\) has an extra symmetry.

\begin{lemma}[window identity]\label{lem:window}
Let \(d_0,\dots,d_{n-1}\) be real numbers with \(d_{n-k}=d_k\) for \(1\le k\le n-1\). Put
\[
P_0=\sum_{\substack{0\le s\le n-1\\ s\text{ even}}}d_s,\qquad
P_1=\sum_{\substack{0\le s\le n-1\\ s\text{ odd}}}d_s.
\]
Then for \(0\le k\le n-1\),
\[
T_{n-1-k}+T_{k-1}=
\begin{cases}
D, & n\text{ odd},\\
2P_1, & n\text{ even and } k\text{ even},\\
2P_0, & n\text{ even and } k\text{ odd}.
\end{cases}
\]
\end{lemma}

\begin{proof}
Let \(g\colon\Z\to\R\) be the \(n\)-periodic function with \(g(s)=d_s\) for
\(0\le s\le n-1\). For \(1\le s\le n-1\) we have \(g(-s)=g(n-s)=d_{n-s}=d_s\) by the
hypothesis. Hence \(g(s)=d_{\abs{s}}\) for all \(\abs{s}\le n-1\), and by \eqref{eq:T},
\(T_m=\sum_{s\in S(m)}g(s)\) for \(0\le m\le n-1\).

Fix \(0\le k\le n-1\). The set \(S(n-1-k)\) consists of the \(n-k\) integers
\(k+1-n+2j\) with \(0\le j\le n-1-k\). If \(k\ge1\), the set \(S(k-1)\) consists of the
\(k\) integers \(1-k+2i\) with \(0\le i\le k-1\). By periodicity
\(g(1-k+2i)=g(n+1-k+2i)\), and \(n+1-k+2i=k+1-n+2j\) for \(j=n-k+i\), which runs through
\(n-k,\dots,n-1\). Therefore (also for \(k=0\), where \(T_{-1}=0\) and \(j\) already runs
through \(0,\dots,n-1\))
\[
T_{n-1-k}+T_{k-1}=\sum_{j=0}^{n-1}g(k+1-n+2j).
\]
If \(n\) is odd, \(2\) is invertible modulo \(n\), so \(j\mapsto k+1-n+2j\) is a
bijection of \(\Z/n\Z\), and the sum is \(\sum_{s=0}^{n-1}g(s)=D\). If \(n\) is even, the
terms for \(j\) and \(j+n/2\) are equal. For \(0\le j\le n/2-1\) the residues of
\(k+1-n+2j\) modulo \(n\) are pairwise distinct, because the numbers \(2j\) are distinct
elements of \(\{0,2,\dots,n-2\}\). They are all congruent to \(k+1\) modulo \(2\), and
exactly \(n/2\) residues modulo \(n\) have that parity, so they are all of them. Hence the
sum is \(2\sum_{0\le s\le n-1,\ s\equiv k+1\ (\mathrm{mod}\ 2)}d_s\), which is \(2P_1\) for even
\(k\) and \(2P_0\) for odd \(k\).
\end{proof}

\begin{lemma}[sector \(A_1\)]\label{lem:A1}
\(\kappa_{\VA}(n)=0\).
\end{lemma}

\begin{proof}
We derive five relations (A.1)--(A.5) from the equations, use Lemma~\ref{lem:window} to
reduce them to a small linear system, and check that this system has only the zero
solution. Let \(v=(r,r,d,d)\in\VA\), with \(r\) palindromic and \(d\) even (coordinates
\(d_0,\dots,d_{n-1}\)), and suppose \((M-nI)v=0\). By \eqref{eq:actA} and
Remark~\ref{rem:components} this means
\begin{align*}
(\mathrm{row}_i)&\quad R+2\sum_{k'=i-n+1}^{i}d_{\abs{k'}}=0\quad(0\le i\le n-1),\\
(\mathrm{diag}_k)&\quad -kd_k+2\rho_k+T_{n-1-k}=0\quad(0\le k\le n-1).
\end{align*}

(A.1) For \(1\le i\le n-1\), subtracting \((\mathrm{row}_{i-1})\) from
\((\mathrm{row}_i)\) gives \(2(d_{\abs{i}}-d_{\abs{i-n}})=2(d_i-d_{n-i})=0\). So
\(d_{n-k}=d_k\) for \(1\le k\le n-1\), and Lemma~\ref{lem:window} applies.

(A.2) \((\mathrm{row}_0)\) reads \(R+2\sum_{k'=-(n-1)}^{0}d_{\abs{k'}}=R+2D=0\).

(A.3) Since \(r\) is palindromic, for \(0\le k\le n\),
\(\rho_{n-k}=\sum_{i\ge n-k}r_i=\sum_{i\ge n-k}r_{n-1-i}=\sum_{j=0}^{k-1}r_j=R-\rho_k\).

(A.4) For \(1\le k\le n-1\), add \((\mathrm{diag}_k)\) and \((\mathrm{diag}_{n-k})\) and
use (A.1) and (A.3):
\begin{align*}
0&=-kd_k-(n-k)d_{n-k}+2(\rho_k+\rho_{n-k})+T_{n-1-k}+T_{k-1}\\
&=-nd_k+2R+T_{n-1-k}+T_{k-1}.
\end{align*}

(A.5) \((\mathrm{diag}_0)\) reads \(2\rho_0+T_{n-1}=2R+T_{n-1}=0\).

\emph{Case \(n\) odd.} By Lemma~\ref{lem:window}, (A.4) becomes \(nd_k=2R+D\) for
\(1\le k\le n-1\), and (A.5) becomes \(2R+D=0\) (Lemma~\ref{lem:window} with \(k=0\)).
Hence \(d_k=0\) for \(1\le k\le n-1\), so \(D=d_0\). Then (A.2) and (A.5) form the system
\[
\begin{pmatrix}1&2\\2&1\end{pmatrix}\begin{pmatrix}R\\D\end{pmatrix}=0,\qquad
\det\begin{pmatrix}1&2\\2&1\end{pmatrix}=1\cdot1-2\cdot2=-3\ne0.
\]
So \(R=D=0\), hence \(d_0=0\) and \(d=0\).

\emph{Case \(n\) even.} By Lemma~\ref{lem:window}, (A.4) becomes \(nd_k=2R+2P_1\) for
even \(k\) and \(nd_k=2R+2P_0\) for odd \(k\) (\(1\le k\le n-1\)), and (A.5) becomes
\(2R+2P_1=0\). Hence \(d_k=0\) for every even \(k\) with \(2\le k\le n-2\), and therefore
\(P_0=d_0\). Also, all \(d_k\) with \(k\) odd equal one number \(\beta\), where
\(n\beta=2R+2d_0\). There are \(n/2\) odd indices in \([1,n-1]\), so
\(P_1=\frac n2\beta\) and \(D=P_0+P_1=d_0+\frac n2\beta\). So \((d_0,\beta,R)\)
satisfies three equations: \(2d_0-n\beta+2R=0\) (the relation for odd \(k\));
\(n\beta+2R=0\) (from (A.5)); and \(2d_0+n\beta+R=0\) (from (A.2)). The second gives
\(R=-\frac n2\beta\). The first then gives \(2d_0=n\beta-2R=2n\beta\). The third then
gives \(2n\beta+n\beta-\frac n2\beta=\frac52n\beta=0\). So \(\beta=0\), \(R=0\),
\(d_0=0\) and \(d=0\). Equivalently, the matrix of this linear system is nonsingular:
\begin{align*}
\det\begin{pmatrix}2&-n&2\\0&n&2\\2&n&1\end{pmatrix}
&=2(n\cdot1-2\cdot n)-(-n)(0\cdot1-2\cdot2)+2(0\cdot n-n\cdot2)\\
&=-2n-4n-4n=-10n\ne0.
\end{align*}

\emph{Conclusion in both cases.} We have \(d=0\) and \(R=0\). Then \((\mathrm{diag}_k)\)
gives \(2\rho_k=kd_k-T_{n-1-k}=0\) for \(0\le k\le n-1\). Together with \(\rho_n=0\) this
gives \(r_i=\rho_i-\rho_{i+1}=0\) for all \(i\). Hence \(v=0\).
\end{proof}

% ==================================================================
\section{The sector \texorpdfstring{\(B_2\)}{B2} at
\texorpdfstring{\(\mu=n\)}{mu=n}}\label{sec:B2}
% ==================================================================

This sector is handled by a discrete maximum principle.

\begin{lemma}[weighted means]\label{lem:means}
Let \(I\) be a finite set, \(j_0\in I\), and \((e_j)_{j\in I}\) real numbers. Suppose that
for every \(j\in I\) with \(j\ne j_0\) there are \(\alpha(j),\beta(j)\in I\) and a real
number \(\lambda_j\) with \(0<\lambda_j\le1\) such that
\[
e_j=\lambda_je_{\alpha(j)}+(1-\lambda_j)e_{\beta(j)},
\]
and that for every \(j\in I\) the sequence \(j,\alpha(j),\alpha(\alpha(j)),\dots\)
reaches \(j_0\) after finitely many steps. Then \(e_j=e_{j_0}\) for all \(j\in I\).
\end{lemma}

\begin{proof}
Let \(m=\max_{j\in I}e_j\) and \(Z=\{j\in I\mid e_j=m\}\), a nonempty set. If \(j\in Z\)
and \(j\ne j_0\), then
\(m=e_j=\lambda_je_{\alpha(j)}+(1-\lambda_j)e_{\beta(j)}\le\lambda_je_{\alpha(j)}
+(1-\lambda_j)m\), so \(\lambda_j(m-e_{\alpha(j)})\le0\). Since \(\lambda_j>0\) and
\(e_{\alpha(j)}\le m\), this forces \(e_{\alpha(j)}=m\), that is, \(\alpha(j)\in Z\).
Take any \(j\in Z\) and follow the sequence \(j,\alpha(j),\alpha(\alpha(j)),\dots\).
Every term before \(j_0\) is reached lies in \(Z\) and differs from \(j_0\), so its
successor lies in \(Z\). Since the sequence reaches \(j_0\), we get \(j_0\in Z\), that
is, \(e_{j_0}=m\). The numbers \((-e_j)_{j\in I}\) satisfy the same hypotheses, so
\(e_{j_0}=\min_je_j\). Hence \(\max_je_j=\min_je_j\), and all \(e_j\) are equal.
\end{proof}

\emph{Setup.} Let \(v=(0,0,d,-d)\in\VBb\) with coordinates \(d_0,\dots,d_{n-1}\). By
\eqref{eq:actB2} and Remark~\ref{rem:components}, \((M-nI)v=0\) if and only if
\begin{equation}\label{eq:Eq}
\Eq(k):\quad kd_k+T_{n-1-k}=0\qquad(0\le k\le n-1).
\end{equation}
Introduce \(e_0=d_0\) and \(e_k=2d_k\) for \(1\le k\le n-1\), with the conventions
\(e_n=e_{n+1}=0\). By the formulas after \eqref{eq:T},
\(T_m=\sum_{0\le j\le m,\ j\equiv m\ (\mathrm{mod}\ 2)}e_j\) for \(0\le m\le n-1\). So
\begin{equation}\label{eq:Eqe}
\Eq(k)=\frac k2e_k+\sum_{\substack{0\le j\le n-1-k\\ j\equiv n-1-k\ (\mathrm{mod}\ 2)}}e_j
\qquad(0\le k\le n-1).
\end{equation}
(For \(k=0\) the first term is \(0\); for \(k\ge1\) it is \(kd_k\).) Define, for
\(0\le k\le n-1\),
\begin{equation}\label{eq:F}
F_k=ke_k-(k+2)e_{k+2}+2e_{n-1-k}.
\end{equation}

\begin{lemma}\label{lem:F}
Put \(\Eq(n)=\Eq(n+1)=0\). Then \(F_k=2\Eq(k)-2\Eq(k+2)\) for \(0\le k\le n-1\), and
\(\Eq(k)=\frac12\sum_{i\ge0,\ k+2i\le n-1}F_{k+2i}\). Hence the systems
\(\{\Eq(k)=0\mid 0\le k\le n-1\}\) and \(\{F_k=0\mid 0\le k\le n-1\}\) have the same
solutions.
\end{lemma}

\begin{proof}
Let \(0\le k\le n-3\). In \eqref{eq:Eqe}, the sums for \(\Eq(k)\) and \(\Eq(k+2)\) run
over the \(j\) of one and the same parity, up to \(n-1-k\) and up to \(n-3-k\)
respectively. So they differ by the single term \(e_{n-1-k}\), and
\(2\Eq(k)-2\Eq(k+2)=ke_k-(k+2)e_{k+2}+2e_{n-1-k}=F_k\). For \(k=n-2\):
\(\Eq(n-2)=\frac{n-2}2e_{n-2}+e_1\) and \(F_{n-2}=(n-2)e_{n-2}-ne_n+2e_1
=(n-2)e_{n-2}+2e_1\). For \(k=n-1\): \(\Eq(n-1)=\frac{n-1}2e_{n-1}+e_0\) and
\(F_{n-1}=(n-1)e_{n-1}-(n+1)e_{n+1}+2e_0=(n-1)e_{n-1}+2e_0\). The second formula follows
by telescoping.
\end{proof}

Two consequences of \eqref{eq:F} are used below:
\begin{gather}
F_0=0\iff e_2=e_{n-1};\label{eq:F0}\\
F_k=0\iff e_{k+2}=\frac{k}{k+2}e_k+\frac{2}{k+2}e_{n-1-k}\qquad(1\le k\le n-3).
\label{eq:Fk}
\end{gather}

\begin{lemma}[sector \(B_2\)]\label{lem:B2}
\begin{enumerate}
\item[(a)] If \(n\) is even, \(\kappa_{\VBb}(n)=0\).
\item[(b)] If \(n\) is odd, \(\kappa_{\VBb}(n)=1\), and \(\ker((M-nI)|_{\VBb})\) is
spanned by \((0,0,\dc,-\dc)\), where \(\dc\) is the even vector with \(\dc_0=n-1\),
\(\dc_k=-1\) for even \(k\) with \(2\le k\le n-1\), and \(\dc_k=0\) for odd \(k\).
\end{enumerate}
\end{lemma}

\begin{proof}
By \eqref{eq:Eq} and Lemma~\ref{lem:F}, \(v\in\ker((M-nI)|_{\VBb})\) if and only if
\(F_k=0\) for \(0\le k\le n-1\). Let \(e\) be such a solution. The relations
\eqref{eq:F0} and \eqref{eq:Fk} express most \(e_j\) as weighted means of two others, so
Lemma~\ref{lem:means} will show that many of them are equal.

(a) \emph{\(n\) even.} Apply Lemma~\ref{lem:means} with \(I=\{1,2,\dots,n-1\}\) and
\(j_0=1\). For \(j=2\): \(e_2=e_{n-1}\) by \eqref{eq:F0}; take
\(\alpha(2)=\beta(2)=n-1\) and \(\lambda_2=1\). For \(3\le j\le n-1\): \eqref{eq:Fk}
with \(k=j-2\in[1,n-3]\) gives \(e_j=\frac{j-2}{j}e_{j-2}+\frac2je_{n+1-j}\); take
\(\alpha(j)=j-2\), \(\beta(j)=n+1-j\in[2,n-2]\subseteq I\) and
\(\lambda_j=\frac{j-2}{j}\in(0,1)\). The sequence \(j,\alpha(j),\dots\) reaches \(1\):
from an odd \(j\) it descends \(j,j-2,\dots,1\); from an even \(j\) it descends
\(j,j-2,\dots,2\), then goes to \(n-1\), which is odd because \(n\) is even, and descends
\(n-1,n-3,\dots,1\). Hence \(e_1=e_2=\dots=e_{n-1}=:c\). Now
\(F_{n-2}=(n-2)e_{n-2}+2e_1=nc\), so \(c=0\); and \(F_{n-1}=(n-1)e_{n-1}+2e_0=2e_0\), so
\(e_0=0\). Thus \(e=0\), so \(d=0\) and \(v=0\).

(b) \emph{\(n\) odd}, so \(n\ge5\). Since \(n-1-k\equiv k\pmod 2\), each \(F_k\) involves
only the \(e_j\) with \(j\equiv k\pmod 2\).

\emph{Even indices.} Apply Lemma~\ref{lem:means} with \(I=\{2,4,\dots,n-1\}\) and
\(j_0=2\). For even \(j\) with \(4\le j\le n-1\), \eqref{eq:Fk} with \(k=j-2\) (even,
\(2\le k\le n-3\)) gives \(e_j=\frac{j-2}je_{j-2}+\frac2je_{n+1-j}\), where \(n+1-j\) is
even and lies in \([2,n-3]\subseteq I\); take \(\alpha(j)=j-2\), \(\beta(j)=n+1-j\) and
\(\lambda_j=\frac{j-2}{j}\in(0,1)\). The sequence \(j,j-2,\dots\) reaches \(2\). Hence
\(e_2=e_4=\dots=e_{n-1}=:c\). Then \(F_{n-1}=(n-1)c+2e_0=0\) gives
\(e_0=-\frac{n-1}2c\).

\emph{Odd indices.} Apply Lemma~\ref{lem:means} with \(I=\{3,5,\dots,n-2\}\) and
\(j_0=3\); for \(n=5\), \(I=\{3\}\) and there is nothing to do. For odd \(j\) with
\(5\le j\le n-2\), \eqref{eq:Fk} with \(k=j-2\) (odd, \(3\le k\le n-4\)) gives
\(e_j=\frac{j-2}je_{j-2}+\frac2je_{n+1-j}\), where \(n+1-j\) is odd and lies in
\([3,n-4]\subseteq I\); take \(\alpha(j)=j-2\), \(\beta(j)=n+1-j\) and
\(\lambda_j=\frac{j-2}j\in(0,1)\). The sequence \(j,j-2,\dots\) reaches \(3\). Hence
\(e_3=e_5=\dots=e_{n-2}=:c'\). Then \(F_{n-2}=(n-2)c'+2e_1=0\) gives
\(e_1=-\frac{n-2}2c'\), and, since the indices \(3\) and \(n-2\) lie in \(I\),
\[
F_1=e_1-3e_3+2e_{n-2}=-\frac{n-2}{2}c'-3c'+2c'=-\frac n2c'.
\]
So \(c'=0\), and \(e_j=0\) for every odd \(j\).

Therefore every solution is \(e=-\frac c2(n-1,0,-2,0,-2,\dots,0,-2)\), that is,
\(d=-\frac c2\dc\). So \(\kappa_{\VBb}(n)\le1\). Conversely, \(\dc\) satisfies
\eqref{eq:Eq}, by direct substitution. For odd \(k\): \(k\dc_k=0\) and
\(T_{n-1-k}=2(\dc_1+\dc_3+\dots+\dc_{n-1-k})=0\), since \(n-1-k\) is odd. For \(k=0\):
\(T_{n-1}=\dc_0+2(\dc_2+\dots+\dc_{n-1})=(n-1)-2\cdot\frac{n-1}2=0\). For even
\(k=2i\ge2\): \(k\dc_k+T_{n-1-k}=-2i+(n-1)-2\bigl(\frac{n-1}2-i\bigr)=0\). Since
\(\dc\ne0\), \(\kappa_{\VBb}(n)=1\).
\end{proof}

Appendix~\ref{app:alt} gives a second, independent proof of Lemma~\ref{lem:B2}. The
Theorem does not depend on it.

% ==================================================================
\section{Proof of the Theorem, and the eigenspace}\label{sec:main}
% ==================================================================

\begin{proof}[Proof of the Theorem]
By Lemma~\ref{lem:gram}, the multiplicity of \(n-4\) is \(\dim\ker(M-nI)\). By
Lemma~\ref{lem:sectors}(d) with \(\mu=n\), and Lemmas~\ref{lem:A1}, \ref{lem:B1},
\ref{lem:B2} and \ref{lem:E} for the four terms,
\[
\dim\ker(M-nI)=\underbrace{0}_{A_1}+\underbrace{\lceil n/2\rceil-1}_{B_1}
+\underbrace{\iv{n\text{ odd}}}_{B_2}+2\cdot\underbrace{0}_{E}.
\]
For even \(n\) this is \(\frac n2-1=\frac{n-2}2\). For odd \(n\) it is
\(\frac{n+1}2-1+1=\frac{n+1}2\). Both numbers are at least \(1\) for \(n\ge4\), so
\(n-4\) is an eigenvalue.
\end{proof}

\begin{corollary}\label{cor:eigenspace}
For \(n\ge4\), the eigenspace of \(\Adj\) for the eigenvalue \(n-4\) consists of the
functions on the board
\[
x(i,j)=r_i-r_j+\gamma\bigl(\iv{i=j}-\iv{i+j=n-1}\bigr),
\]
where \(r\in\R^n\) is palindromic and \(\gamma\in\R\), with \(\gamma=0\) if \(n\) is even.
This description is due to Cardoso, Costa and Duarte; see Remark~\ref{rem:credit}.
\end{corollary}

\begin{proof}
By Lemma~\ref{lem:gram}, the eigenspace is \(N\ker(M-nI)\). By the proof of
Lemma~\ref{lem:sectors}(d), \(\ker(M-nI)\) is the direct sum of the kernels of \(M-nI\)
on the five sectors. The kernels on \(\VA\), \(\VE\) and \(\VEp\) are zero
(Lemmas~\ref{lem:A1} and \ref{lem:E}, and Lemma~\ref{lem:sectors}(c) for \(\VEp\)). So,
by Lemmas~\ref{lem:B1} and \ref{lem:B2}, \(\ker(M-nI)\) is spanned by the vectors
\((r,-r,0,0)\) with \(r\) palindromic and \(\sum_ir_i=0\), together with
\((0,0,\dc,-\dc)\) when \(n\) is odd. By \eqref{eq:Nv}, \(N(r,-r,0,0)(i,j)=r_i-r_j\).
This expression does not change when a constant is added to \(r\), so the condition
\(\sum_ir_i=0\) can be dropped. For odd \(n\),
\(N(0,0,\dc,-\dc)(i,j)=\dc_{\abs{i-j}}-\dc_{\abs{i+j-(n-1)}}\). Here \(i-j\) and
\(i+j-(n-1)\) have the same parity, because \(n-1\) is even. If both are odd, the value
is \(0-0=0\). If both are even, the value is \((n-1)-(-1)=n\) when \(i=j\) and
\(i+j\ne n-1\); it is \((-1)-(n-1)=-n\) when \(i\ne j\) and \(i+j=n-1\); and it is \(0\)
when both or neither of \(i=j\), \(i+j=n-1\) hold. So
\(N(0,0,\dc,-\dc)=n\bigl(\iv{i=j}-\iv{i+j=n-1}\bigr)\).
\end{proof}

\begin{remark}[credit for the eigenspace]\label{rem:credit}
Cardoso, Costa and Duarte proved that the multiplicity of \(n-4\) is at least the value
in the Theorem, by exhibiting linearly independent eigenvectors. They did so first in
\cite[Theorem~3.10]{CCD20}, for \(n\ge3\), and restated it in
\cite[Theorem~4.1 and Corollary~4.2]{CCD23}. With squares indexed from \(1\), their
eigenvectors are of two kinds. The first kind (the set \(Y^n\) defined after
\cite[Definition~3.9]{CCD20}, the vectors \(F_{n,\ell}\) of \cite{CCD23}) is the
indicator of the columns \(\ell\) and \(n+1-\ell\) minus the indicator of the rows
\(\ell\) and \(n+1-\ell\). (When \(n\) is odd and \(\ell=(n+1)/2\), the two columns
coincide, as do the two rows; the vector of \(Y^n\) is then twice this indicator
difference, while \(F_{n,\ell}\) equals it.) In our notation
this is \(x(i,j)=r_i-r_j\) with \(r=-\iv{i\in\{\ell-1,n-\ell\}}\), a palindromic \(r\).
So the first summand in Corollary~\ref{cor:eigenspace} is the span of these vectors. The
second kind, used for odd \(n\) (the set \(Z^n\) of \cite{CCD20}, the vector
\(E_{n,n-4}\) of \cite{CCD23}), is a nonzero multiple of the indicator of the main
diagonal minus the indicator of the antidiagonal (also defined after
\cite[Definition~3.9]{CCD20}); with our indexing it is a multiple of
\(\iv{i=j}-\iv{i+j=n-1}\), the \(\gamma\)-term. As printed in \cite{CCD20}, the set \(Y^n\) has
\(\lceil n/2\rceil\) vectors with one linear relation, and the proof there uses
\(\lceil n/2\rceil-1\) of them, together with \(Z^n\) for odd \(n\), as linearly independent
eigenvectors; so the number of linearly independent ones, \((n-2)/2\) resp.\ \((n+1)/2\), is a
lower bound on the multiplicity. Hence their eigenvectors span the space
described in Corollary~\ref{cor:eigenspace}, and once the Theorem bounds the
multiplicity from above, they span the whole eigenspace. The description of the
eigenspace is therefore theirs.
The exact value, equivalently the upper bound (the Theorem), is conjectured but not proved in their papers (\cite[Conjecture~3]{CCD20}; \cite[Conjecture~1]{CCD23}; \cite[Conjecture~4.1]{CCD25}); whether a proof of it,
or the method used here, appears elsewhere has not been checked beyond our literature search of 2026-11-27.
We checked these theorem numbers and definitions against the PDFs of both papers on
2026-11-27, and the identification by exact computation for small \(n\).
\end{remark}

% ==================================================================
\section{Computer verification}\label{sec:computer}
% ==================================================================

No step of the proof uses a computer. The computations in this section confirm the
Theorem, and the individual steps of its proof, for finitely many \(n\). All of them use
exact integer or rational arithmetic; none uses floating point. We call a computation a
\emph{certificate} when it has been rerun from a clean copy, reproduced by an
independent second implementation, and archived with its tool versions. None of the
checks below has yet been archived as a certificate in that sense, so we label them
\emph{numerical}.

The scripts and their archived outputs are in the folder \nolinkurl{scripts/} of the
accompanying repository,
\begin{center}
\footnotesize\url{https://github.com/nebulousice/cardoso-costa-duarte-queens-conjecture},
\end{center}
and every command below is run from that folder. The software was
Python~3.12.10, python-flint~0.7.1 and sympy~1.14.0, on Windows~11. The file
\nolinkurl{REPRODUCE.md} gives runtimes and the expected output lines.

\begin{center}
\small
\begin{tabular}{@{}>{\raggedright\arraybackslash}p{0.40\textwidth}%
>{\raggedright\arraybackslash}p{0.17\textwidth}%
>{\raggedright\arraybackslash}p{0.31\textwidth}@{}}
\toprule
Check & Range & Status\\
\midrule
Nullity of \(M-nI\), full matrix, two rank algorithms & \(4\le n\le40\) &
numerical (exact arithmetic, two implementations)\\
Nullity of \(M-nI\) on each sector & \(4\le n\le150\) & numerical (exact arithmetic)\\
Replay of every identity in the proof & \(4\le n\le40\) & numerical (exact arithmetic)\\
Nullity of \(\Adj-(n-4)I\) & \(4\le n\le12\) & numerical (exact arithmetic)\\
Individual steps of Lemmas~\ref{lem:sectors}, \ref{lem:action}, \ref{lem:E},
\ref{lem:A1}, \ref{lem:B2}, the appendix and Lemma~G & \(4\le n\le40\);
\(q\le80\) & numerical (exact arithmetic)\\
\bottomrule
\end{tabular}
\end{center}

\begin{enumerate}
\item \emph{Nullity of \(M-nI\), \(4\le n\le40\).} For each \(n\) the integer matrix
\(M-nI\) is built square by square from the board, with no symmetry used. Its rank over
\(\mathbb{Q}\) is computed twice: by FLINT's exact integer rank, and by pure-Python
fraction-free (Bareiss) elimination. Both runs build the matrix with the same routine;
they differ in the rank algorithm. For every \(n\) in the range both give nullity
\((n-2)/2\) for even \(n\) and \((n+1)/2\) for odd \(n\). Commands:
\begin{verbatim}
python p2_verify.py fullflint 4 40
python p2_verify.py fraction 4 16
python p2_verify.py fraction 17 40
\end{verbatim}
Archived outputs: \nolinkurl{scripts/logs/fullflint_4_40.log},
\nolinkurl{scripts/logs/fraction_4_16.log}, \nolinkurl{scripts/logs/fraction_17_40.log}.

The same script also checks three groups of proof steps in exact rational arithmetic,
for every \(n\) with \(4\le n\le40\): the reduced linear systems in the proof of
Lemma~\ref{lem:A1}, both parts of the appendix, and Lemma~G for \(3\le q\le40\).
Command:
\begin{verbatim}
python p2_verify.py lemmas 4 40
\end{verbatim}
Archived output: \nolinkurl{scripts/logs/verify_lemmas_4_40.log}.

\item \emph{Nullities on the sectors.} The sector matrices are written from
\eqref{eq:actA}--\eqref{eq:actE}, and their nullities at \(\mu=n\) are computed with
FLINT. They are \(0\) for \(A_1\), \(\lceil n/2\rceil-1\) for \(B_1\),
\(\iv{n\text{ odd}}\) for \(B_2\) and \(0\) for \(E\), for every \(n\) with
\(4\le n\le150\). For \(41\le n\le150\) this is the only computation. Commands:
\begin{verbatim}
python p2_verify.py sectors 4 40
python p2_verify.py sectors 41 150
\end{verbatim}
The archived outputs are the files \nolinkurl{sectors_flint_4_40.log} and
\nolinkurl{sectors_flint_41_150.log} in the folder \nolinkurl{scripts/logs/}.

\item \emph{Replay of every identity in the proof, \(4\le n\le40\).} A separate
pure-Python script rebuilds everything from the board. For each \(n\) it checks
Lemma~\ref{lem:clique} (for \(n\le24\)), the relations of Lemma~\ref{lem:sym}, the
sector decomposition, the formulas of Lemma~\ref{lem:action}, the identities in the
proofs of Lemmas~\ref{lem:E}, \ref{lem:window} and \ref{lem:A1}, Lemma~\ref{lem:F} and
the conclusions of Lemma~\ref{lem:B2}, and \(\dim\ker(M-nI)\). The last check uses the
rank modulo the prime \(2^{31}-1\), which bounds the nullity from above, together with
explicit kernel vectors, which bound it from below. It also checks
Corollary~\ref{cor:eigenspace} on the adjacency matrix (\(n\le16\)) and the appendix
(\(q\le60\)). Command:
\begin{verbatim}
python prover_check.py 4 40
\end{verbatim}
Archived output: \nolinkurl{scripts/logs/replay_4_40.log}. Its last line is
\texttt{PROVER REPLAY: ALL PASS}.

\item \emph{The adjacency matrix directly, \(4\le n\le12\).} The nullity of
\(\Adj-(n-4)I\) is computed on the \(n^2\times n^2\) matrix with FLINT, and the vectors
of Corollary~\ref{cor:eigenspace} are checked to lie in the kernel and to have full
rank. Command:
\begin{verbatim}
python p2_verify.py adjacency 4 12
\end{verbatim}
Archived output: \nolinkurl{scripts/logs/adjacency_4_12.log}.

\item \emph{Individual steps.} A third script, written independently of the other two,
checks from the board the sector formulas, the steps of Lemmas~\ref{lem:E},
\ref{lem:A1} and \ref{lem:B2}, and the pivots in the proof of Lemma~G
(\(1\le q\le80\)). Commands:
\begin{verbatim}
python p2_proofcheck.py sectors 4 40
python p2_proofcheck.py a1 4 40
python p2_proofcheck.py e 4 40
python p2_proofcheck.py b2 4 40
python p2_proofcheck.py g 1 80
python p2_proofcheck.py symbolic
\end{verbatim}
Archived outputs: the files \nolinkurl{scripts/logs/proofcheck_*.log}.
\end{enumerate}

% ==================================================================
\section{Formal verification}\label{sec:formal}
% ==================================================================

The Theorem and Corollary~\ref{cor:eigenspace} are formalized in the proof assistant
Lean~4 (version 4.37.0-rc1), using the
mathematics library Mathlib at the commit pinned in the file
\nolinkurl{lean/lake-manifest.json} of the accompanying repository. ``Formalized'' means
that Lean checks a complete proof of the statement, step by step, from the definitions,
and that no step is assumed without proof. This section says what the Lean files state,
what a reader has to trust, how the check was run, how the Lean proof is organized, and
what is not formalized. The Lean files are in the folder
\nolinkurl{lean/Pipeline/QueensSpectrum/} of the accompanying repository (12 files,
2283 lines); the file \nolinkurl{LEAN.md} at the top level of the repository gives more
detail.

\subsection{The definition of the graph}
In Lean a square is a pair \(x=(x_1,x_2)\) of numbers in \(\{0,1,\dots,n-1\}\), as in
Section~\ref{sec:notation}. The file \nolinkurl{Defs.lean} defines the relation
\(\mathrm{Attack}\) (lines~27--29) and the graph \lean{queensGraph n} (adjacency
rule on line~34). In mathematical notation they read
\begin{align*}
\mathrm{Attack}(x,y)&\iff x_1=y_1\ \text{or}\ x_2=y_2\ \text{or}\ x_1-x_2=y_1-y_2\\
&\qquad\text{or}\ x_1+x_2=y_1+y_2,\\
x\text{ and }y\text{ are adjacent}&\iff x\ne y\ \text{and}\ \mathrm{Attack}(x,y).
\end{align*}
The differences and sums are computed in \(\Z\), so there is no arithmetic modulo \(n\).
This is the definition of \(\QG(n)\) in Section~\ref{sec:notation}. These four lines of
Lean are the only definitions of this work that appear in the statements of the Theorem
and of Corollary~\ref{cor:eigenspace}.

\subsection{The statement proved}
The main results are in the file \nolinkurl{Main.lean}. Let \(K\) be a \emph{linearly
ordered field}: a field with a total order compatible with addition and multiplication,
such as \(\mathbb{Q}\) or \(\R\). Let \(\Adj_K\) be the adjacency matrix of \(\QG(n)\)
with entries in \(K\) (Mathlib's \lean{SimpleGraph.adjMatrix}). For every such
\(K\) and every integer \(n\ge4\), Lean proves the following.
\begin{itemize}
\item \lean{queens\_mult}: the dimension over \(K\) of
\(\ker\bigl(\Adj_K-(n-4)I\bigr)\) is \((n-2)/2\) if \(n\) is even and \((n+1)/2\) if \(n\)
is odd.
\item \lean{queens\_alg\_mult}: the multiplicity of \(n-4\) as a root of the
characteristic polynomial of \(\Adj_K\) is the same number.
\item \lean{queens\_hasEigenvalue}: \(n-4\) is an eigenvalue of \(\Adj_K\).
\end{itemize}
The theorems \lean{queens\_mult\_real}, \lean{queens\_mult\_rat} and
\lean{queens\_alg\_mult\_real} are the cases \(K=\R\) and \(K=\mathbb{Q}\). In the
Lean text the right-hand side is a natural number computed with truncating division.
For \(n\ge4\) the division is exact, so the value is the stated one.

For \(K=\R\) this is exactly the Theorem of Section~\ref{sec:intro}, with both readings of
``multiplicity'' (the dimension of the kernel, and the order as a root of the
characteristic polynomial). The source indexes squares from \(1\) and Lean from \(0\);
shifting the indices is a graph isomorphism.

Corollary~\ref{cor:eigenspace} is in the file \nolinkurl{Eigenspace.lean}. For every
linearly ordered field \(K\), every \(n\ge4\) and every function \(x\) from the squares
to \(K\), Lean proves the following.
\begin{itemize}
\item \lean{queens\_eigenspace}: \(x\) lies in the eigenspace of \(\Adj_K\) for the
eigenvalue \(n-4\) if and only if there are \(r\in K^n\) with \(r_{n-1-i}=r_i\) for all
\(i\), and \(\gamma\in K\) with \(\gamma=0\) when \(n\) is even, such that
\(x(i,j)=r_i-r_j+\gamma\bigl(\iv{i=j}-\iv{i+j=n-1}\bigr)\) for all squares \((i,j)\).
\end{itemize}
The theorem \lean{queens\_eigenspace\_real} is the case \(K=\R\). Both inclusions are
proved, so this is exactly Corollary~\ref{cor:eigenspace}.

\subsection{What has to be trusted}
A reader who accepts the Lean check needs to trust only the following.
\begin{enumerate}
\item The Lean kernel, which checks every proof.
\item The four lines of definition shown above.
\item Mathlib's standard notions used in the statement: the adjacency matrix of a
simple graph, the linear map given by a matrix, the eigenspace of a linear map, the
dimension of a vector space, the characteristic polynomial, the multiplicity of a root,
the reversal \(i\mapsto n-1-i\) of \(\{0,\dots,n-1\}\) (used to state that \(r\) is
palindromic in Corollary~\ref{cor:eigenspace}), and the notions of field, linear order
and even number.
\item The three standard axioms of Lean's logic:
\begin{center}
\lean{propext},\qquad \lean{Classical.choice},\qquad \lean{Quot.sound}.
\end{center}
The command \lean{\#print axioms} reports that the main theorems use these three and
nothing else.
\end{enumerate}
The files contain no unproved step (no \lean{sorry} or \lean{admit}). They
do not use \lean{native\_decide}, which would add trust in compiled code. The other
definitions in the files (the lines, the incidence matrix, the three symmetries, the
sectors) are used only inside the proof. A mistake in one of them could make the proof
fail to check; it could not make a false statement pass.

\subsection{How it was checked}
The check is one command, run from the root of the accompanying repository:
\begin{verbatim}
bash scripts/lean_check.sh QueensSpectrum
\end{verbatim}
It builds the project with \lean{lake build}, scans the files for
\lean{sorry} and \lean{admit}, and prints the axioms of each main theorem.
It ends with a line beginning \lean{PASS lean\_check QueensSpectrum} only if the
build has no error, the scan is empty and only the three standard axioms occur. On
2026-11-27 it passed in two separate runs, and again in two separate runs after the file
\nolinkurl{Eigenspace.lean} for Corollary~\ref{cor:eigenspace} was added. A further AI
referee (Claude Opus) rebuilt the files of the Theorem from a fresh copy, again with a
pass and identical axiom output.

The same referee checked that the Lean statement matches the source conjecture. It also
compared the Lean definition with an independently written queen-move rule. Lean's
kernel confirmed that the two agree on all \(256\) pairs of squares for \(n=4\), and
evaluation confirmed agreement for every \(n\le8\). The number of edges equals the
source's formula \(n(n-1)(5n-1)/3\) for \(n\le8\). For \(n=4,5,6\) the Lean adjacency
matrices agree entry by entry with matrices built by an independent script. In a second
check, the same model confirmed, clause by clause, that the Lean statement of
\lean{queens\_eigenspace} says what Corollary~\ref{cor:eigenspace} says. Its own exact
rational computation for \(4\le n\le10\) found that the functions of the Lean statement
span the kernel of \(\Adj-(n-4)I\), and that the eigenvectors printed by Cardoso, Costa
and Duarte span the same space.

\emph{Kernel replays.} A Lean build checks each proof with Lean's kernel, the small
program that accepts or rejects proofs, while the file is compiled. A kernel replay checks
the compiled proofs again afterwards, from the stored files. On 2026-11-27, Lean's own
checker \lean{leanchecker}, shipped with Lean 4.37.0-rc1, replayed every declaration of
all 12 files, and its separately built variant \lean{leanchecker-paranoid} did the same.
With the option \lean{-{}-fresh}, both also re-added \nolinkurl{Main.lean} and everything
it imports, Lean's core library and Mathlib included, to an empty environment. All these
runs passed. Two kernels written independently of Lean's own, \lean{nanoda\_bin}
(version 0.4.31) and \lean{lean4lean}, checked an export file of the nine main theorems
and every declaration they depend on, made with the exporter \lean{lean4export}
(version 3.1.0). Both passed, as did two further kernels shipped with the toolchain. A
copy of the export file in which one proof was replaced by the proof of another theorem
was rejected by all five checkers that read export files. One limit remains: the independent kernels
check the exporter's rendering of the compiled files, not the compiled files themselves.
Only Lean's own kernel replayed those directly. The records are in the folder
\nolinkurl{lean/kernel-replays/} of the repository.

The file \nolinkurl{REPRODUCE.md} gives versions, runtimes and the expected output.

\subsection{How the Lean proof is organized}
The Lean proof follows the written proof. The table lists where each step is
formalized, by file and Lean name.
\begin{center}
\small
\begin{tabular}{@{}>{\raggedright\arraybackslash}p{0.30\textwidth}%
>{\raggedright\arraybackslash}p{0.62\textwidth}@{}}
\toprule
Written proof & Lean (file: names)\\
\midrule
Section~\ref{sec:notation}, \eqref{eq:Nv} & \nolinkurl{Defs}: \lean{Line}, \lean{N}, \lean{M}, \lean{N\_mulVec}\\
Lemma~\ref{lem:clique} & \nolinkurl{Defs}: \lean{lemma1}\\
Lemma~\ref{lem:gram} & \nolinkurl{Reduction}: \lean{lemma2}\\
Lemma~\ref{lem:sym} & \nolinkurl{Symmetry}: \lean{lt\_lh}, \lean{lh\_lt\_lh}, \lean{M\_mulVec\_comp} and others\\
Lemma~\ref{lem:sectors}(a)--(c) & \nolinkurl{Symmetry}, \nolinkurl{Kernel}: \lean{decomp}, \lean{W\_zero}, the projectors\\
Lemma~\ref{lem:action}, Remark~\ref{rem:components} & \nolinkurl{Defs}: \lean{M\_mulVec\_row}, \lean{M\_mulVec\_diag}; \nolinkurl{Symmetry}: \lean{ker\_eqs}\\
Lemma~\ref{lem:B1} & \nolinkurl{Kernel}: \lean{B1\_form}, \lean{vecB1\_mem}; \nolinkurl{Count}: \lean{finrank\_P0}\\
Lemma~\ref{lem:E} & \nolinkurl{SectorE}: \lean{sectorE\_pure}; \nolinkurl{Kernel}: \lean{E\_zero}\\
Lemmas~\ref{lem:window} and \ref{lem:A1} & \nolinkurl{SectorA1}: \lean{sectorA1\_pure}; \nolinkurl{Kernel}: \lean{A1\_zero}\\
Lemma~\ref{lem:means} & \nolinkurl{SectorB2}: \lean{weighted\_means}\\
Lemmas~\ref{lem:F} and \ref{lem:B2} & \nolinkurl{SectorB2}: \lean{sectorB2\_pure}; \nolinkurl{Kernel}: \lean{B2\_form}, \lean{vecB2\_mem}\\
Proof of the Theorem & \nolinkurl{Kernel}: \lean{ker\_char}; \nolinkurl{Count}: \lean{finrank\_E}; \nolinkurl{Main}: \lean{queens\_mult}\\
Algebraic multiplicity & \nolinkurl{AlgMult}: the two multiplicities agree for symmetric matrices\\
Corollary~\ref{cor:eigenspace} & \nolinkurl{Eigenspace}: \lean{N\_vecB1}, \lean{N\_vecB2}, \lean{queens\_eigenspace}\\
\bottomrule
\end{tabular}
\end{center}
The comments in the Lean files letter the three parts of Lemma~\ref{lem:sectors}
differently: their parts (a), (b), (c) are parts (b), (c), (a) of
Lemma~\ref{lem:sectors} here.

\subsection{Where the Lean proof differs from the written proof}
The Lean proof reaches the same conclusions, sometimes by a different route. No
statement is weakened. The differences are these.
\begin{enumerate}
\item \emph{The count.} Lemma~\ref{lem:sectors}(d) is not formalized as stated. Instead,
Lean shows directly that \(\ker(M-nI)\) consists exactly of the vectors
\((r,-r,0,0)+\gamma(0,0,\dc,-\dc)\) with \(r\) palindromic, \(\sum_ir_i=0\), and
\(\gamma=0\) for even \(n\). Here \(\dc\) is the vector of Lemma~\ref{lem:B2}(b). Lean
then computes the dimension through an injective linear parametrization. The number is
the same as in Section~\ref{sec:main}.
\item \emph{Lemma~\ref{lem:window}.} Lean uses a weaker consequence of the window
identity, which suffices: the number \(U_k=T_{n-1-k}+T_{k-1}\) satisfies \(U_{k+2}=U_k\). This
follows from \(T_{m+2}=T_m+2d_{m+2}\) and the symmetry \(d_{n-k}=d_k\). The small linear
systems at the end of the proof of Lemma~\ref{lem:A1} are solved by explicit linear
combinations, so their determinants are not computed as such.
\item \emph{Lemma~\ref{lem:means}.} Lean assumes a rank function on the index set that
strictly decreases along \(\alpha\) away from \(j_0\), instead of assuming that the chain
\(j,\alpha(j),\alpha(\alpha(j)),\dots\) reaches \(j_0\). For the index maps used in
Lemma~\ref{lem:B2} the two hypotheses are equivalent.
\item \emph{Lemma~\ref{lem:E}.} The last step uses an algebraic identity instead of a
positivity argument (the positive bracket in Step~3). The equations \((\mathrm{diag}_k)\) and Steps 1 and 2 of its proof give
\(kd_k=\rho_k=\rho_{n-k}=(n-k)d_{n-k}=-(n-k)d_k\), hence \(nd_k=0\).
\item \emph{Lemma~\ref{lem:F}.} Only the direction used in the proof is formalized:
the equations \(\Eq\) imply the equations \(F\).
\item \emph{Dimensions.} The dimensions of the sectors in
Lemma~\ref{lem:sectors}(a), and the check after Lemma~\ref{lem:sectors}, are not
formalized. The proof does not need them.
\end{enumerate}

\subsection{What is not formalized}
\begin{itemize}
\item The appendix (a second proof of Lemma~\ref{lem:B2}). The Theorem does not use it.
\item Remark~\ref{rem:credit} and the other statements about the literature. They are
about published papers, which Lean cannot check.
\item The computer checks of Section~\ref{sec:computer}. For the Theorem they are
superseded by the Lean proof. They remain independent checks of the definitions.
\item Parts (a) and (c) of the conjecture. This paper does not address them.
\end{itemize}

% ==================================================================
\appendix
\section{A second proof of \texorpdfstring{Lemma~\ref{lem:B2}}{Lemma 12}}\label{app:alt}
% ==================================================================

This appendix is not used by the Theorem. It gives independent proofs of both parts of
Lemma~\ref{lem:B2}. Notation is as in Section~\ref{sec:B2}.

\subsection{Even \texorpdfstring{\(n\)}{n}: an explicit recurrence}
Let \(n=2p\). Put \(E_j=T_{2j}\) for \(0\le j\le p-1\), and \(Q_j=T_{2(p-1-j)+1}\) for
\(0\le j\le p-1\), with \(Q_p=0\). Since \(T_0=d_0\), \(T_1=2d_1\) and
\(T_m-T_{m-2}=2d_m\) for \(2\le m\le n-1\), the numbers \(E_j\) and \(Q_j\) determine
\(d\): \(d_0=E_0\), \(2d_{2i}=E_i-E_{i-1}\) for \(1\le i\le p-1\), and
\(2d_{2i+1}=Q_{p-1-i}-Q_{p-i}\) for \(0\le i\le p-1\). The equations \eqref{eq:Eq}
become:
\begin{itemize}
\item \(\Eq(0)\): \(T_{2p-1}=Q_0=0\);
\item \(\Eq(2i)\), \(1\le i\le p-1\): \(i(E_i-E_{i-1})+Q_i=0\), that is,
\(E_i=E_{i-1}-Q_i/i\);
\item \(\Eq(2i+1)\), \(0\le i\le p-1\), written with \(j=p-1-i\):
\(\frac{2p-1-2j}{2}(Q_j-Q_{j+1})+E_j=0\), that is,
\(Q_{j+1}=Q_j+\frac{2E_j}{2p-1-2j}\).
\end{itemize}
By induction on \(j\), starting from \(Q_0=0\),
\[
E_j=E_0\,\frac{2p-1-2j}{2p-1}\quad(0\le j\le p-1),\qquad
Q_j=E_0\,\frac{2j}{2p-1}\quad(0\le j\le p).
\]
Indeed, \(Q_{j+1}=E_0\bigl(\frac{2j}{2p-1}+\frac{2}{2p-1}\bigr)\) and
\(E_{j+1}=E_0\bigl(\frac{2p-1-2j}{2p-1}-\frac{2}{2p-1}\bigr)\). The convention
\(Q_p=0\) now reads \(E_0\frac{2p}{2p-1}=0\). So \(E_0=0\), all \(E_j\) and \(Q_j\)
vanish, and \(d=0\). This proves Lemma~\ref{lem:B2}(a) again.

\subsection{Odd \texorpdfstring{\(n\)}{n}: a block form}
Let \(n=2p+1\) with \(p\ge2\). Put \(u_l=e_{2l}\) for \(0\le l\le p\) and
\(w_l=e_{2l+1}\) for \(0\le l\le p-1\). For \(q\ge0\) let \(G_q\) be the symmetric matrix
of size \(q+1\), with rows and columns indexed by \(0,\dots,q\), and entries
\((G_q)_{ab}=a\iv{a=b}+\iv{a+b\le q}\). By \eqref{eq:Eqe}, for \(0\le i\le p\) and
\(0\le i\le p-1\) respectively,
\begin{align*}
\Eq(2i)&=iu_i+\sum_{l=0}^{p-i}u_l=(G_pu)_i,\\
\Eq(2i+1)&=\frac{2i+1}{2}w_i+\sum_{l=0}^{p-1-i}w_l
=\tfrac12\bigl((I+2G_{p-1})w\bigr)_i.
\end{align*}

\begin{lemmaG}
For \(q\ge1\), the matrix \(G_q\) is positive semidefinite of rank \(q\), and its kernel
is spanned by \((q,-1,-1,\dots,-1)\).
\end{lemmaG}

\begin{proof}
We eliminate the indices \(q,q-1,\dots,1\) one at a time and track the pivots. We use
the following fact. Let \(X\) be a real symmetric matrix indexed by a finite set \(J\),
let \(k\in J\) with \(X_{kk}\ne0\), and let \(X'\) be the matrix indexed by
\(J\setminus\{k\}\) with \(X'_{ab}=X_{ab}-X_{ak}X_{kb}/X_{kk}\). Then \(X\) is congruent
to the block-diagonal matrix with blocks \(X'\) and \((X_{kk})\). Hence
\(\rank X=\rank X'+1\), and \(X\) is positive semidefinite if and only if \(X_{kk}>0\)
and \(X'\) is positive semidefinite.

Let \(m=\lfloor q/2\rfloor\). Write \(X^{(k)}\) for the matrix on \(\{0,\dots,k\}\)
obtained after eliminating \(q,\dots,k+1\); so \(X^{(q)}=G_q\).

\emph{Step 1: the indices \(q,\dots,m+1\).} We claim that for \(m\le k\le q\) and
\(0\le a,b\le k\),
\[
X^{(k)}_{ab}=(G_q)_{ab}-\sum_{i=k+1}^{q}\frac{\iv{a\le q-i}\,\iv{b\le q-i}}{i}.
\]
This holds for \(k=q\). Suppose it holds for some \(k\ge m+1\). For every \(i\ge k+1\)
we have \(k+i\ge2m+3>q\), so \(\iv{k\le q-i}=0\), and row \(k\) of \(X^{(k)}\) equals
row \(k\) of \(G_q\). Thus \(X^{(k)}_{kk}=k+\iv{2k\le q}=k\), as \(2k\ge2m+2>q\), and
\(X^{(k)}_{ak}=\iv{a\le q-k}\) for \(a<k\). So the pivot is \(k>0\), and eliminating
\(k\) subtracts \(\iv{a\le q-k}\,\iv{b\le q-k}/k\). This gives the claim for \(k-1\). For
\(k=m\) and \(0\le a,b\le m\) we have \(a+b\le2m\le q\), so
\((G_q)_{ab}=a\iv{a=b}+1\), and the claim reads
\[
X^{(m)}_{ab}=a\iv{a=b}+c_{\max(a,b)},\qquad
c_t=1-\sum_{i=m+1}^{q-t}\frac1i\quad(0\le t\le m),
\]
where the sum is empty when \(q-t<m+1\) (this happens only for \(q=2m\), \(t=m\)).

\emph{Step 2: the indices \(m,\dots,1\).} Put \(\gamma_k=\frac{k}{q-k}\) for
\(0\le k\le m\); note \(q-k\ge q-m\ge1\). We claim that for \(0\le k\le m\) and
\(0\le a,b\le k\),
\[
X^{(k)}_{ab}=a\iv{a=b}+c_{\max(a,b)}-c_k+\gamma_k.
\]
For \(k=m\) this is Step~1, provided \(c_m=\gamma_m\). If \(q=2m\), then
\(c_m=1=\frac{m}{q-m}\). If \(q=2m+1\), then \(c_m=1-\frac1{m+1}=\frac m{m+1}
=\frac{m}{q-m}\). Suppose the claim holds for some \(k\) with \(1\le k\le m\). Then
\(X^{(k)}_{kb}=\gamma_k\) for \(b<k\), and the pivot is
\[
\pi_k=X^{(k)}_{kk}=k+\gamma_k=\frac{k(q+1-k)}{q-k}>0.
\]
Eliminating \(k\) subtracts the constant \(\gamma_k^2/\pi_k\) from every entry. Now
\(\gamma_k-\gamma_k^2/\pi_k=\gamma_k(\pi_k-\gamma_k)/\pi_k=k\gamma_k/\pi_k
=\frac{k}{q+1-k}\). Also \(c_{k-1}-c_k=-\frac1{q+1-k}\): the sum defining \(c_{k-1}\)
has one extra term, \(i=q+1-k\), which is at least \(m+1\) because \(k\le m\le q-m\).
Hence
\begin{align*}
-c_k+\frac{k}{q+1-k}&=-c_{k-1}-\frac{1}{q+1-k}+\frac{k}{q+1-k}\\
&=-c_{k-1}+\frac{k-1}{q-(k-1)}=-c_{k-1}+\gamma_{k-1},
\end{align*}
which is the claim for \(k-1\).

\emph{Conclusion.} For \(k=0\) the claim gives the \(1\times1\) matrix
\(X^{(0)}=(c_0-c_0+\gamma_0)=(0)\). By the fact, applied \(q\) times, \(G_q\) is
congruent to the diagonal matrix with entries \(q,q-1,\dots,m+1,\pi_m,\dots,\pi_1,0\),
of which the first \(q\) are positive. So \(G_q\) is positive semidefinite of rank
\(q\). Finally \(\kappa=(q,-1,\dots,-1)\) lies in the kernel: \((G_q\kappa)_0=q-q=0\),
and for \(1\le a\le q\), \((G_q\kappa)_a=-a+\bigl(q-(q-a)\bigr)=0\), because row \(a\)
of the matrix \(\bigl(\iv{a+b\le q}\bigr)_{a,b}\) has ones exactly in the columns
\(0,\dots,q-a\).
\end{proof}

\begin{proof}[Second proof of Lemma~\ref{lem:B2}(b)]
The system \eqref{eq:Eq} is \(G_pu=0\) together with \((I+2G_{p-1})w=0\). By Lemma~G
with \(q=p\ge2\), \(\ker G_p\) is spanned by \(u=(p,-1,\dots,-1)\); that is,
\(d_0=p\) and \(d_{2l}=-\frac12\) for \(1\le l\le p\), which is \(\frac12\dc\) on the
even indices. By Lemma~G with \(q=p-1\ge1\), \(G_{p-1}\) is positive semidefinite, so
\(I+2G_{p-1}\) is positive definite and \(w=0\). Hence the solutions of \eqref{eq:Eq}
form a space of dimension exactly \(1\).
\end{proof}

% ==================================================================
\section*{Acknowledgements}
% ==================================================================

We thank Domingos M.~Cardoso, In\^es Ser\^odio Costa and Rui Duarte for the conjecture,
the lower bounds and the eigenvectors on which this paper builds.

The three Claude referees named in the abstract reviewed an earlier text of the same
proof. The present text is a re-presentation of that refereed text, with no mathematical
change. ChatGPT and Gemini read the first version of the present text. The second version
added the Lean formalization of the Theorem. This third version adds the Lean
formalization of Corollary~\ref{cor:eigenspace} and the kernel replays
(Section~\ref{sec:formal}), and corrects several citations and bibliographic details. The
mathematical argument is unchanged.

% ==================================================================
\begin{thebibliography}{BCHV16}
% ==================================================================

\bibitem[BCHV16]{BCHV16}
A.~E. Brouwer, S.~M. Cioab\u{a}, W.~H. Haemers and J.~R. Vermette,
\emph{Notes on simplicial rook graphs},
J. Algebraic Combin. \textbf{43} (2016), no.~4, 783--799,
\href{https://doi.org/10.1007/s10801-015-0633-y}{doi:10.1007/s10801-015-0633-y};
\href{https://arxiv.org/abs/1408.5615}{arXiv:1408.5615}.

\bibitem[CCD20]{CCD20}
D.~M. Cardoso, I.~S. Costa and R.~Duarte,
\emph{Spectral properties of the \(n\)-Queens' graphs}, preprint, 2020;
\href{https://arxiv.org/abs/2012.01992}{arXiv:2012.01992}.

\bibitem[CCD22]{CCD22}
D.~M. Cardoso, I.~S. Costa and R.~Duarte,
\emph{Sharp bounds on the least eigenvalue of a graph determined from edge clique
partitions}, J. Algebraic Combin. \textbf{58} (2023), no.~1, 263--277,
\href{https://doi.org/10.1007/s10801-023-01247-1}{doi:10.1007/s10801-023-01247-1};
\href{https://arxiv.org/abs/2201.01224}{arXiv:2201.01224}.

\bibitem[CCD23]{CCD23}
D.~M. Cardoso, I.~S. Costa and R.~Duarte,
\emph{Integer eigenvalues of the \(n\)-Queens graph},
Discrete Math. \textbf{347} (2024), no.~9, Paper No.~114065,
\href{https://doi.org/10.1016/j.disc.2024.114065}{doi:10.1016/j.disc.2024.114065};
\href{https://arxiv.org/abs/2301.08106}{arXiv:2301.08106}.

\bibitem[CCD25]{CCD25}
D.~M. Cardoso, I.~S. Costa and R.~Duarte,
\emph{A family of regular integral graphs and its application to the \(n\)-Queens'
graph}, preprint, 2025;
\href{https://arxiv.org/abs/2508.01088}{arXiv:2508.01088}.

\bibitem[Cos24]{Cos24}
I.~F.~S. Costa, \emph{The \(n\)-Queens graph and its generalizations}, PhD thesis,
Universidade de Aveiro, 2024;
\href{http://hdl.handle.net/10773/41899}{hdl:10773/41899}.

\bibitem[MW15]{MW15}
J.~L. Martin and J.~D. Wagner,
\emph{On the spectra of simplicial rook graphs}, Graphs Combin. \textbf{31} (2015),
no.~5, 1589--1611,
\href{https://doi.org/10.1007/s00373-014-1452-y}{doi:10.1007/s00373-014-1452-y};
\href{https://arxiv.org/abs/1209.3493}{arXiv:1209.3493}.

\end{thebibliography}

\end{document}
